% polar_bias/master/master.tex -- master theorem document. Build: pdflatex master && bibtex master && pdflatex master (x2)
\documentclass[10pt]{article}
\usepackage[utf8]{inputenc}
\usepackage[T1]{fontenc}
\usepackage{amsmath,amssymb,amsthm,mathtools}
\usepackage{booktabs,longtable,array}
\usepackage{enumitem}
\usepackage[margin=1in]{geometry}
\usepackage{xcolor}
\usepackage[round]{natbib}
\usepackage[colorlinks=true,linkcolor=blue!55!black,citecolor=blue!55!black,urlcolor=blue!55!black]{hyperref}
\setlength{\emergencystretch}{2em}
\newtheorem{theorem}{Theorem}
\newtheorem{proposition}[theorem]{Proposition}
\newtheorem{lemma}[theorem]{Lemma}
\newtheorem{corollary}[theorem]{Corollary}
\newtheorem{conjecture}[theorem]{Conjecture}
\theoremstyle{definition}
\newtheorem{definition}[theorem]{Definition}
\newtheorem{assumption}{Assumption}
\renewcommand{\theassumption}{\Alph{assumption}}
\theoremstyle{remark}
\newtheorem{remark}[theorem]{Remark}
% Shared math macros (master document and both papers). No theorem environments here.
\providecommand{\R}{\mathbb{R}}
\providecommand{\Z}{\mathbb{Z}}
\providecommand{\E}{\mathbb{E}}
\providecommand{\Prob}{\mathbb{P}}
\providecommand{\nnuc}[1]{\|#1\|_{*}}
\providecommand{\nF}[1]{\|#1\|_{F}}
\providecommand{\inner}[2]{\langle #1, #2\rangle}
\providecommand{\tr}{\operatorname{tr}}
\providecommand{\polar}{\operatorname{polar}}
\providecommand{\Med}{\operatorname{Med}}
\providecommand{\dist}{\operatorname{dist}}
\providecommand{\vol}{\operatorname{vol}}
\providecommand{\rank}{\operatorname{rank}}
\providecommand{\Cov}{\operatorname{Cov}}
\providecommand{\He}{\mathrm{He}}
\providecommand{\cA}{\mathsf{A}}
\providecommand{\cB}{\mathsf{B}}
\providecommand{\ones}{\mathbf{1}}
\providecommand{\st}{\dagger}
\providecommand{\ew}{\varepsilon_w}
\providecommand{\Lw}{L_w}
\providecommand{\SO}{\mathrm{SO}}
\providecommand{\Id}{\mathrm{Id}}
\providecommand{\gs}{g_\ast}
\providecommand{\beff}{b_{\mathrm{eff}}}
\providecommand{\cK}{\mathcal{K}}
\providecommand{\cL}{\mathcal{L}}
\providecommand{\Le}{L_\varepsilon}
\providecommand{\psie}{p}
\providecommand{\argmin}{\operatorname*{arg\,min}}

% Shared statements (bodies only). Each document wraps them in its own theorem environments and labels.
% ---------------------------------------------------------------------------------------------
% Definitions and assumptions
% ---------------------------------------------------------------------------------------------
\newcommand{\StmtClass}{%
Fix constants $0<\lambda\le\Lambda$, $C_\ast\ge1$, $c_\ast\in(0,1]$. For $\varepsilon\in(0,1]$, a random vector $X$ in $E$
belongs to the class $\mathcal X(\varepsilon)$ if there are a covariance $\Sigma$ with
$\lambda\,\Id\preceq\Sigma\preceq\Lambda\,\Id$, a symmetric $3$-tensor $\cK$ with $|\cK|\le C_\ast\varepsilon$, and an even
polynomial $\widehat r$ of degree at most $6$ with $\widehat r(0)=0$ and coefficients bounded by $C_\ast\varepsilon^2$,
such that, with the \emph{comparison transform}
$\widehat\psie(t)=e^{-\inner t{\Sigma t}/2}\big(1+\tfrac16\cK[(it)^{\otimes3}]+\widehat r(t)\big)$ and
$\chi_X(t)=\E e^{i\inner tX}$, $a=c_\ast/\varepsilon$:
\begin{itemize}
\item[\textup{(E1)}] $\int_{t_0+[-a,a]^N}|\chi_X(t)|\,dt\le C_\ast$ for every $t_0\in E$;
\item[\textup{(E2)}] $\int_{[-a,a]^N}|\chi_X(t)-\widehat\psie(t)|\,dt\le C_\ast\varepsilon^3$;
\item[\textup{(E3)}] $\Prob(|X|>u)\le C_\ast e^{-u^2/C_\ast}$ for all $u\ge0$.
\end{itemize}
The \emph{comparison density} is the signed density $\psie$ with $\int e^{i\inner ty}\psie(y)dy=\widehat\psie(t)$; explicitly
$\psie=\varphi_\Sigma\,(1+\frac16\inner{\cK}{\He_3^\Sigma}+r)$ with $r$ an even polynomial. Constants denoted $C,c$ in
statements about $\mathcal X(\varepsilon)$ depend only on $N$, $\nu$, $\lambda,\Lambda,C_\ast,c_\ast$.}

\newcommand{\StmtAssB}{%
$\zeta,\zeta_1,\zeta_2,\dots$ are i.i.d.\ in $E$ with $|\zeta|\le M$ almost surely, $\E\zeta=0$, and covariance
$\Sigma\succeq\lambda\,\Id$, $\lambda>0$; $K_3=\E\zeta^{\otimes3}$. A \emph{weight sequence} is $w=(w_k)_{k\ge1}$ with
$w_k\ge0$, $\sum_kw_k=1$; we set
\[
\rho_p=\sum_kw_k^p,\quad\ \xi_w=\sum_kw_k\zeta_k,\quad\ X_w=\rho_2^{-1/2}\xi_w,\quad\
\ew=\frac{\max_kw_k}{\sqrt{\rho_2}},\quad\ \Lw=\log(e/\ew).
\]
No condition is imposed on the characteristic function of $\zeta$.}

\newcommand{\StmtAssR}{%
There is a closed cone $V\subset E$ such that $o\in C^\infty(E\setminus V)$,
$|D^ko(y)|\le C_k\dist(y,V)^{-k}$ for all $k\ge0$, and
$\vol\{y\in B(0,1):\dist(y,V)\le r\}\le C_Vr^{c}$ for $r\le1$.}

\newcommand{\StmtAssL}{%
$N\ge3$, and there are $\kappa\in(0,1]$ with $2+\kappa<N$ and a constant $C$ such that the Fourier transform of the
finite signed measure $\lambda_\nu=\Delta(\chi\nu)$ satisfies $|\widehat{\lambda_\nu}(t)|\le C(1+|t|)^{-\kappa}$ for all
$t\in E$. Here $\chi$ is a fixed smooth radial cut-off equal to $1$ on $B(0,1)$ and supported in $B(0,2)$.}

\newcommand{\StmtAssN}{%
$\xi_b=\frac1b\sum_{i\le b}\zeta_i$ with $\zeta_i$ i.i.d.\ in $E$, $\E\zeta=0$, $\E|\zeta|^3<\infty$, covariance
$\Sigma\succ0$, third cumulant tensor $K_3=\E\,\zeta^{\otimes3}$, and Cram\'er's condition
$\limsup_{|t|\to\infty}|\E e^{i\inner t\zeta}|<1$.}

% ---------------------------------------------------------------------------------------------
% Paper II statements
% ---------------------------------------------------------------------------------------------
\newcommand{\StmtCodim}{%
Let Assumption~\ref{ass:R} hold with $c\ge2$. There are $\varepsilon_0,C>0$ such that for every
$X\in\mathcal X(\varepsilon)$ with $\varepsilon\le\varepsilon_0$, writing $\Le=\log(e/\varepsilon)$,
\begin{equation}\label{eq:codimX}
\argmin_{x\in E}\ \E\big[\nu(x+X)-\nu(X)\big]\ \subset\ B\big(\hat x,\ C\varepsilon^{\min(c,3)}\Le^{N/2}\big),\qquad
\hat x=-\tfrac16J_\Sigma^{-1}T_\Sigma[\cK],
\end{equation}
In particular, under Assumption~\ref{ass:B}, for every weight sequence with $\ew\le\varepsilon_0$,
\begin{equation}\label{eq:codimw}
\arg\min\psi_{\xi_w}\ \subset\ B\Big(\frac{\rho_3}{\rho_2}\,\gs\,,\ C\sqrt{\rho_2}\;\ew^{\min(c,3)}\,\Lw^{N/2}\Big),
\end{equation}
so that for the mean of $b$ samples
$\Med_\nu(\xi_b)=-\gs/b+O(b^{-3/2}\log^{N/2}b)$ if $c=2$ and $\Med_\nu(\xi_b)=-\gs/b+O(b^{-2}\log^{N/2}b)$ if $c\ge3$.}

\newcommand{\StmtLaplace}{%
Let Assumption~\ref{ass:L} hold and put $q=N/(N-\kappa)\in(1,\frac32]$. There are $\varepsilon_0,C>0$ such that for every
$X\in\mathcal X(\varepsilon)$ with $\varepsilon\le\varepsilon_0$, with $H(x)=\E[\nu(x+X)-\nu(X)]$,
$\Psi(x)=\int(\nu(x+y)-\nu(y))\psie(y)dy$ and $d_\varepsilon=\varepsilon^{q}\Le^{(N+1)/4}$:
\begin{equation}\label{eq:laplaceX}
\sup_{|x|\le2}|H(x)-\Psi(x)|\le Cd_\varepsilon^2,\qquad
\sup_{|x|\le1}\ \sup_{v\in\partial H(x)}|v-\nabla\Psi(x)|\le Cd_\varepsilon,\qquad
\arg\min H\subset B(\hat x,Cd_\varepsilon),
\end{equation}
with $\hat x=-\frac16J_\Sigma^{-1}T_\Sigma[\cK]$. In particular, under Assumption~\ref{ass:B}, for every weight sequence
with $\ew\le\varepsilon_0$,
\begin{equation}\label{eq:laplacew}
\arg\min\psi_{\xi_w}\ \subset\ B\Big(\frac{\rho_3}{\rho_2}\,\gs\,,\ C\sqrt{\rho_2}\;\ew^{q}\,\Lw^{(N+1)/4}\Big),
\end{equation}
so that $\Med_\nu(\xi_b)=-\gs/b+O\big(b^{-1-\frac{\kappa}{2(N-\kappa)}}\log^{\frac{N+1}4}b\big)$ for the mean of $b$ samples.
The middle estimate of \eqref{eq:laplaceX} holds for \emph{every} subgradient, hence for the expectation of every
measurable selection of $\partial\nu(x+X)$, whatever its values on the singular set.}

\newcommand{\StmtNuclear}{%
Let $\nu=\nnuc\cdot$ on $\R^{K\times d}$ with $(K,d)\ne(1,1)$, $N=Kd$, and let Assumption~\ref{ass:B} hold. Then for every
bounded law with nondegenerate covariance, lattice laws included,
\[
\Med_\ast(\xi_b)=-\frac{\gs}b+O\big(b^{-1-\gamma}\log^{N}b\big),\qquad
\gamma=\begin{cases}1&\text{if }|K-d|\ge2,\\[1pt] \tfrac12&\text{if }|K-d|=1,\\[1pt]
\tfrac1{14}&\text{if }K=d=2,\\[1pt] \tfrac1{4K^2-2}&\text{if }K=d\ge3,\end{cases}
\]
and more generally $\arg\min\psi_{\xi_w}\subset B\big(\frac{\rho_3}{\rho_2}\gs,\,C\,(\max_kw_k)\,\ew^{2\gamma}\Lw^{N}\big)$,
uniformly over weight sequences with $\ew\le\varepsilon_0$. For $K=d=1$ the
statement is false (Theorem~\ref{thm:sign}).}

\newcommand{\StmtSign}{%
Let $\zeta$ be a real random variable with $\E\zeta=0$, variance $\sigma^2>0$, $\E|\zeta|^3<\infty$,
$\kappa_3=\E\zeta^3$, and $\gs=\kappa_3/(6\sigma^2)$; let $S_b=\zeta_1+\dots+\zeta_b$ and
$G_b^\st=\arg\min_G\E|G+S_b/b|$.
\begin{enumerate}
\item[\textup{(a)}] If $\zeta$ is not lattice, $b\,G_b^\st\to\gs$ (every selection).
\item[\textup{(b)}] If $\zeta$ is lattice, supported on $a+h\Z$ with maximal span $h$, let $\Lambda_b=-(ba+h\Z)$. For every
$\delta>0$ and $b\ge b_0(\delta)$: $b\,G^\st_b\subset[\gs-\frac h2-\delta,\ \gs+\frac h2+\delta]$, and if a point
$\ell$ of $\Lambda_b$ satisfies $|\ell-\gs|\le\frac h2-\delta$ then $b\,G^\st_b=\{\ell\}$. Consequently:
\begin{enumerate}
\item[\textup{(i)}] if $a/h$ is irrational, the limit points of $b\,G_b^\st$ fill the interval $[\gs-\frac h2,\gs+\frac h2]$;
\item[\textup{(ii)}] if $a/h=p/q$ in lowest terms, then for $b$ in each residue class $j$ modulo $q$ for which $\gs$ is not
a midpoint of $-(\frac{jh}q+h\Z)$, $b\,G_b^\st$ is eventually equal to the point of that coset nearest to $\gs$; there
is no limit if $q\ge2$, and if $q=1$ and $\gs\notin\frac h2+h\Z$ the limit is $h\operatorname{round}(\gs/h)$;
\item[\textup{(iii)}] if $\zeta\in h\Z$ and $|\kappa_3|<3\sigma^2h$, then $G_b^\st=\{0\}$ for all large $b$: sign descent is
exactly unbiased, although $\gs\ne0$ when $\kappa_3\ne0$.
\end{enumerate}
\end{enumerate}
In the lattice case $b\,G_b^\st\to\gs$ only if $a\in h\Z$ and $\gs\in h\Z$.}

\newcommand{\StmtWor}{%
Let $n\ge3$ and $z_1,\dots,z_n\in E$ with $\sum_iz_i=0$, $|z_i|\le M$, $\Sigma_n=\frac1n\sum_iz_iz_i^\top\succeq\lambda\,\Id$,
$K_{3,n}=\frac1n\sum_iz_i^{\otimes3}$, and let $\gs^{(n)}=-\frac16J_{\Sigma_n}^{-1}T_{\Sigma_n}[K_{3,n}]$. Let
$\xi_{n,b}$ be the mean of a uniformly random subset of size $b\in\{1,\dots,n-1\}$, $p=b/n$ and $v=np(1-p)=b(n-b)/n$. Suppose that $\nu$
satisfies Assumption~\ref{ass:R} with $c\ge2$, and put $\gamma=\frac12(\min(c,3)-1)$, or Assumption~\ref{ass:L}, and put
$\gamma=\frac{\kappa}{2(N-\kappa)}$. Then there are $v_0,C$, depending only on $\nu,M,\lambda$, such that whenever
$v\ge v_0$,
\begin{equation}\label{eq:wor}
\sup_{m\in\Med_\nu(\xi_{n,b})}\Big|\,b\,m+\frac{n-2b}{n-2}\,\gs^{(n)}\Big|\le C\,v^{-\gamma}\,\log^{N}(ev).
\end{equation}
In particular the first-order bias of a batch drawn without replacement is that of a batch drawn with replacement
multiplied by $\frac{n-2b}{n-2}$: it vanishes at $b=n/2$ and \emph{changes sign} beyond, for every bounded population
with nondegenerate covariance, lattice populations included, as soon as $b\to\infty$ and $n-b\to\infty$.}

\newcommand{\StmtMomCor}{%
Let $N_\beta=(1-\beta)\sum_{k\ge0}\beta^k\xi_b^{(k)}$ with independent batch means $\xi_b^{(k)}$ of samples satisfying
Assumption~\ref{ass:B}, and let $\beff=b\,\frac{1+\beta+\beta^2}{1-\beta^2}$, $\theta=\frac{1-\beta^2}b$. If $\nu$ satisfies
Assumption~\ref{ass:R} with $c\ge2$ (put $\gamma=\frac12(\min(c,3)-1)$) or Assumption~\ref{ass:L} (put
$\gamma=\frac\kappa{2(N-\kappa)}$), then
\[
\Med_\nu(N_\beta)=-\frac{\gs}{\beff}+O\Big(\frac{\theta^{\gamma}\log^{N}(e/\theta)}{\beff}\Big)
\qquad\text{uniformly in }b\ge1,\ \beta\in[0,1)\text{ with }\theta\le\varepsilon_0^2 .
\]
In particular the first-order law holds at \emph{fixed} batch size, even $b=1$, as $\beta\to1$.}

% ---------------------------------------------------------------------------------------------
% Paper I statements
% ---------------------------------------------------------------------------------------------
\newcommand{\StmtConv}{%
Let $F:E\to\R$ be differentiable, strongly convex, with Lipschitz gradient; let $\nu$ be any norm on $E$ and
$(\xi_t)_{t\ge1}$ i.i.d.\ with $\E|\xi_t|^2<\infty$. From deterministic $W_1,M_0$ run
\[
M_t=\beta M_{t-1}+(1-\beta)\big(\nabla F(W_t)+\xi_t\big),\qquad W_{t+1}=W_t-\eta_tO_t,\qquad O_t\in\partial\nu(M_t),
\]
with $\beta\in[0,1)$ fixed, $O_t$ any measurable function of $(\xi_1,\dots,\xi_t)$ with values in $\partial\nu(M_t)$,
and $\eta_t=\eta_0(t+t_0)^{-\theta}$, $\eta_0,t_0>0$, $\frac12<\theta\le1$. Let
$N_\beta=(1-\beta)\sum_{j\ge0}\beta^j\xi_j$ be the stationary filtered noise. Then, almost surely, $\nabla F(W_t)$ converges
to a (possibly random) point
\[
G_\infty\in\arg\min\psi_{N_\beta}=-\Med_\nu(N_\beta),\qquad\text{and}\qquad W_t\to(\nabla F)^{-1}(G_\infty).
\]
No assumption is made on the law of the noise beyond a second moment: it may be discrete, lattice, or carried by a
proper subspace, and $\partial\nu(M_t)$ may be multivalued with positive probability at every step.}

\newcommand{\StmtCF}{%
Let $\nu$ be the $\ell_1$, the Euclidean or the nuclear norm with its direction map, and let Assumption~\ref{ass:N}
hold. Then for every $\epsilon>0$ there is $b_0$ such that for all $b\ge b_0$
\begin{equation}\label{eq:cf}
\arg\min\psi_{\xi_b}\ \subset\ B\Big(\frac{\gs}b,\frac\epsilon b\Big),\qquad\text{i.e.}\qquad
\Med_\nu(\xi_b)=-\frac{\gs}b+o\Big(\frac1b\Big),\qquad \gs=-\tfrac16\,J_\Sigma^{-1}\,T_\Sigma[K_3].
\end{equation}}

\newcommand{\StmtTwo}{%
Let $K,d\ge2$, $r=\min(K,d)$, $\bar d=\max(K,d)$, $\nu=\nnuc\cdot$ on $\R^{K\times d}$, and let $Z$ be a standard Gaussian
$K\times d$ matrix with singular values $s_1,\dots,s_r$. Put
\begin{equation}\label{eq:q}
m_p=\E\sum_is_i^p,\qquad q_{K,d}=\frac1{m_1}\,\E\sum_{i\ne j}\frac{s_is_j}{s_i+s_j},\qquad c_{K,d}=\frac{m_1}{Kd},
\end{equation}
\begin{equation}\label{eq:a12}
a_1=\frac{3\big[(K-1)(d-1)-2(K+d+1)\,q_{K,d}\big]}{(K-1)(d-1)(K+2)(d+2)},\qquad
a_2=\frac{6\big[(Kd+2)\,q_{K,d}+(K-1)(d-1)\big]}{(K-1)(d-1)(K+2)(d+2)} .
\end{equation}
\begin{enumerate}
\item[\textup{(a)}] \emph{(Reduction.)} For every symmetric $3$-tensor $K_3$ and $\sigma>0$,
$J_{\sigma^2I}=\sigma^{-1}c_{K,d}\,\Id$ and $T_{\sigma^2I}[K_3]=-\sigma^{-3}c_{K,d}\,(a_1\cA[K_3]+a_2\cB[K_3])$, where
$\cA[K_3]_{ij}=\sum_{k,l}K_3^{(ij),(kl),(kl)}$ and $\cB[K_3]_{ij}=\sum_{k,l}K_3^{(kj),(kl),(il)}$. Hence, for a noise law
with covariance $\sigma^2I$ and third cumulant $K_3=\E\zeta^{\otimes3}$,
\begin{equation}\label{eq:two}
\gs=\frac1{6\sigma^2}\Big(a_1\,\E\big[\zeta\nF\zeta^2\big]+a_2\,\E\big[\zeta\zeta^\top\zeta\big]\Big).
\end{equation}
The bias operator $K_3\mapsto\gs$ has rank $Kd$ on a space of dimension $\binom{Kd+2}3$.
\item[\textup{(b)}] \emph{(One number.)} $m_3=(K+d-q_{K,d})\,m_1$; $0<q_{K,d}<\frac{r-1}2$; $a_2>0$.
\item[\textup{(c)}] \emph{(Rank-one and full-rank directions.)}
$a_1+a_2=\dfrac{3(3+2q_{K,d})}{(K+2)(d+2)}$ and $r\,a_1+a_2=\dfrac{3(\bar d-1-2q_{K,d})}{(\bar d-1)(\bar d+2)}>0$.
\item[\textup{(d)}] \emph{(Negative fourth derivative.)} For every $D\ne0$,
$\nabla^4\psi_I(0)[D,D,D,D]=-c_{K,d}\big(a_1\nF D^4+a_2\tr(DD^\top)^2\big)<0$.
\item[\textup{(e)}] \emph{(Proportional limit.)} If $K,d\to\infty$ with $K/d\to\gamma\in(0,1]$, then $q_{K,d}/K\to\kappa(\gamma)$,
$d\,a_2\to6\kappa(\gamma)$ and $d^2a_1\to\frac3\gamma\big(1-2(1+\gamma)\kappa(\gamma)\big)$, where
$\kappa(\gamma)=\big[(1+\gamma)\mu_{1/2}(\gamma)-\mu_{3/2}(\gamma)\big]/\big(\gamma\,\mu_{1/2}(\gamma)\big)$ and
$\mu_s(\gamma)=\int x^s\,d\mathrm{MP}_\gamma(x)$ are the moments of the Marchenko--Pastur law; $\kappa(1)=\frac25$.
\end{enumerate}
For $K=1$ the two contractions coincide and $\gs=\frac{3}{d+2}\cdot\frac{\E[\zeta|\zeta|^2]}{6\sigma^2}$.}

\newcommand{\StmtAone}{%
For all $K,d\ge2$,
\[
q_{K,d}>\frac{(K-1)(d-1)}{2(K+d+1)},\qquad\text{equivalently}\qquad a_1<0<a_2,\qquad\text{equivalently}\qquad
a_1+a_2>\frac3{K+d+1}.
\]
With $W=ZZ^\top$, $D_p=\E\tr W^{p/2}$ and $r(p)=D_{p+2}/D_p$, each of the following is also equivalent:
\begin{enumerate}
\item[\textup{(i)}] $2\,r(1)+r(4)<3\,r(2)$: a three-point concavity of the moment ratio of a real Wishart matrix;
\item[\textup{(ii)}] $p\mapsto\E\sum_j s_j^pA_j/\E\sum_js_j^p$, with $A_j=\sum_{i\ne j}s_i^2$, is strictly convex at the
points $p=1,2,4$;
\item[\textup{(iii)}] for the size-biased mean spectral measure $\varpi(dx)\propto x\,\E\sum_i\delta_{s_i^2}(dx)$, of mean
$m=K+d+1$ and variance $v=Kd+K+d+3$: $\ \E_\varpi[x^{1/2}]/\E_\varpi[x^{-1/2}]<m-\frac v{2m}$.
\end{enumerate}
For $K=1$, and in \textup{(iii)} for every Gamma law, there is equality.}

\newcommand{\StmtDim}{%
For independent entries of variance $\sigma^2$ and third cumulant $\kappa_3$, $\cA=\cB=\kappa_3\ones\ones^\top$ and
\[
\gs=\frac{\kappa_3}{6\sigma^2}\,(a_1+a_2)\,\ones\ones^\top,\qquad
a_1+a_2=\frac{3\,(3+2q_{K,d})}{(K+2)(d+2)},\qquad
\frac{3}{K+d+1}<a_1+a_2<\frac{3}{\max(K,d)+2}\quad(K,d\ge2).
\]
Per entry, the first-order offset of a polar step is the scalar Cornish--Fisher offset $\kappa_3/(6\sigma^2b)$ of sign
descent divided by a factor between $\frac13(\max(K,d)+2)$ and $\frac13(K+d+1)$. In the proportional limit
$\max(K,d)(a_1+a_2)\to6\kappa(\gamma)$, equal to $\frac{12}5$ for square matrices.}

\newcommand{\StmtRankOne}{%
If $\zeta$ has rank one almost surely (per-sample gradients $\delta x^\top$ of a linear layer) and covariance
$\sigma^2I$, then $\zeta\zeta^\top\zeta=\nF\zeta^2\zeta$ and, for the nuclear norm,
$\gs=\frac{3+2q_{K,d}}{(K+2)(d+2)}\cdot\frac{\E[\zeta\nF\zeta^2]}{2\sigma^2}$, whereas for the Frobenius norm (normalized
steps) $\gs=\frac{1}{Kd+2}\cdot\frac{\E[\zeta\nF\zeta^2]}{2\sigma^2}$. Orthogonalization multiplies the first-order offset
of normalization by
\[
\varrho_{K,d}=\frac{(3+2q_{K,d})(Kd+2)}{(K+2)(d+2)},\qquad \frac{Kd+2}{K+d+1}<\varrho_{K,d}<\frac{Kd+2}{\max(K,d)+2}<\min(K,d).
\]}

\newcommand{\StmtCompare}{%
\textup{(i)} To first order a polar step sees the third cumulant only through $a_1\cA+a_2\cB$ ($Kd$ numbers), a
normalized step only through $\cA$, and sign descent only through the $Kd$ own third cumulants
$K_3^{(ij),(ij),(ij)}$.
\textup{(ii)} \emph{Sign-blind noise biases polar steps:} if $\zeta$ has covariance $\sigma^2I$ and all third cumulants zero
except $\E[\zeta_{11}\zeta_{12}^2]=\tau\ne0$ (and its permutations), sign descent has no first-order offset and
$\gs=\frac{(a_1+a_2)\tau}{6\sigma^2}e_1e_1^\top$.
\textup{(iii)} \emph{Polar-blind noise biases sign descent:} the tensor supported on $\{(11),(22)\}$ with
$K_3^{(11)^3}=1$, $K_3^{(11),(22),(22)}=-\frac{a_1+a_2}{a_1}$, $K_3^{(22)^3}=K_3^{(22),(11),(11)}=0$ has
$a_1\cA+a_2\cB=0$ and a nonzero own cumulant.
\textup{(iv)} \emph{One skewed factor:} if $\zeta=\epsilon M+\zeta'$ with $\epsilon$ a scalar of third cumulant
$\kappa_3(\epsilon)$, $M$ a fixed matrix, $\zeta'$ independent and symmetric, and total covariance $\sigma^2I$, then
$\gs=\frac{\kappa_3(\epsilon)}{6\sigma^2}(a_1\nF M^2M+a_2MM^\top M)$ and $\inner{\gs}{\kappa_3(\epsilon)M}>0$.}

\newcommand{\StmtBeff}{%
Let $\nu$ and the samples be as in Theorem~\ref{thm:cf}, and let $N_w=\sum_{k\ge0}w_k\xi_b^{(k)}$ for independent batch
means $\xi_b^{(k)}$ and fixed real weights with $\sum_k|w_k|<\infty$, $w\ne0$, $\rho_p=\sum_kw_k^p$. Then
\[
\arg\min\psi_{N_w}\subset B\Big(\frac{\rho_3}{\rho_2}\cdot\frac{\gs}b,\ o\Big(\frac1b\Big)\Big)\qquad(b\to\infty).
\]
For the exponential moving average $w_k=(1-\beta)\beta^k$ this is $\gs/\beff+o(1/b)$ with the \emph{skewness-effective
batch size}
\[
\beff=b\,\frac{1+\beta+\beta^2}{1-\beta^2},\qquad\text{not the variance-effective one}\quad b\,\frac{1+\beta}{1-\beta}.
\]}

\newcommand{\StmtState}{%
Let $\hat g_b(W)=\frac1b\sum_{i\le b}g(W;\omega_i)$ with $\nabla F(W)=\E g(W;\omega)$, let $\Sigma(W)$, $K_3(W)$ be the
covariance and third cumulant of $g(W;\omega)-\nabla F(W)$, and $h_b(W)=\E\,o(\hat g_b(W))$ the mean field. Assume, on a
bounded neighbourhood $\mathcal U$ of $W_\ast$: \textup{(i)} $F$ is $C^1$, twice differentiable at $W_\ast$,
$\nabla F(W_\ast)=0$, $H=\nabla^2F(W_\ast)$ is invertible and $|\nabla F(W)|\ge c|W-W_\ast|$; \textup{(ii)} $\Sigma(W)$,
$K_3(W)$ are bounded, continuous at $W_\ast$, and $\Sigma(W)\succeq\lambda I$; \textup{(iii)} the one-term expansion of
Lemma~\ref{lem:edge} holds uniformly in $W\in\mathcal U$. Let $\gs$ be computed from $\Sigma(W_\ast)$, $K_3(W_\ast)$.
Then for every $\epsilon>0$ there are $b_0$ and $\delta>0$ such that for $b\ge b_0$ every $W\in\mathcal U$ with
$|h_b(W)|\le\delta b^{-1/2}$, in particular every zero of the mean field in $\mathcal U$, satisfies
\[
\Big|W-W_\ast-\frac1b\,H^{-1}\gs\Big|\le\frac\epsilon b .
\]
If moreover $h_b$ is continuous on $\mathcal U$, it has a zero in $\mathcal U$ for $b$ large.}

\newcommand{\StmtLin}{%
Suppose $g(W;\omega)-\nabla F(W)=L_W\chi(\omega)$ on $\mathcal U$, where $\chi\in\R^p$ satisfies Assumption~\ref{ass:N}
and $W\mapsto L_W$ is continuous with $L_W:\R^p\to E$ onto. Then \textup{(ii)} and \textup{(iii)} of
Theorem~\ref{thm:state} hold on a neighbourhood of $W_\ast$.}

\newcommand{\StmtFS}{%
Let $F=\frac1n\sum_if_i$ on $\R^{K\times d}$, $(K,d)\ne(1,1)$, and at a point $W$ let $\zeta=\nabla f_I(W)-\nabla F(W)$ with
$I$ uniform on $\{1,\dots,n\}$. If the $n$ gradients $\nabla f_i(W)$ affinely span $\R^{K\times d}$ (so $n\ge Kd+1$),
then for minibatches of size $b$ drawn with replacement the nuclear medians of the batch noise are
$-\gs/b+O(b^{-1-\gamma}\log^{Kd}b)$ with $\gamma$ as in Theorem~\ref{thm:nuclear}, whatever the arithmetic of the
gradients. When $n=Kd+1$ the noise is a lattice law for \emph{every} data set: no weak Cram\'er or non-lattice
condition can hold, and the law holds nevertheless.}

\newcommand{\StmtLattice}{%
Let $E=\R^2$ with the Euclidean norm, and let $\zeta$ be supported on $a+\cL$ for a lattice $\cL$ of covolume $|\cL|$
generated by the support of $\zeta-\zeta'$. Let $\Lambda_b=ba+\cL$ and
\[
\mathfrak E_\Lambda(m)=\lim_{\tau\to\infty}\Big[|\cL|\sum_{\ell\in\Lambda}\frac{m+\ell}{|m+\ell|}e^{-\inner\ell{\Sigma^{-1}\ell}/2\tau^2}
-\int_{\R^2}\frac{m+y}{|m+y|}e^{-\inner y{\Sigma^{-1}y}/2\tau^2}dy\Big].
\]
Then $b\,G_b^\st=\gs-b^{-1/2}\,(2\pi\sqrt{\det\Sigma})^{-1}J_\Sigma^{-1}\mathfrak E_{\Lambda_b}(\gs)+o(b^{-1/2})$: the
second-order term of $G_b^\st$ is of exact order $b^{-3/2}$, and depends on $b$ only through the coset $\Lambda_b$. If so, the remainder of Theorem~\ref{thm:codim} is
attained for $c=2$.}

\newcommand{\status}[1]{\textsf{\small[#1]}}

\title{\vspace{-2em}The bias law of norm-steepest stochastic steps\\[2pt]
\large Master document: statements, assumptions, dependencies, proofs and verification record}
\author{Project ``Accelerating Variance Reduction'', folder \texttt{polar\_bias/master}}
\date{25 November 2026}
\begin{document}
\maketitle
\tableofcontents

\section{Purpose and conventions}

This document consolidates, in one notation and with complete proofs, the mathematics of three research notes
(\texttt{polar\_bias/paper/polar\_bias}, \texttt{polar\_bias/note2/paper/discrete\_bias}, and the note \texttt{orbital-bias}) and the
results of the original manuscript \texttt{paper/paper\_saga\_to\_muon} that they build on. The proofs are the files
\texttt{proofs/p\_*.tex}; the two papers in \texttt{polar\_bias/papers} input the same files, so that a statement and its proof
exist in exactly one place. Section~\ref{sec:deps} lists every result with its hypotheses, what it depends on, and how
it was verified; Section~\ref{sec:verif} is the verification record.

\paragraph{Status tags.} \status{P} proved here in full. \status{P*} proved here, using a classical theorem that is quoted
(named at the place of use). \status{CA} proved with a computer-assisted step in exact arithmetic. \status{F} formal
derivation supported by exact numerics; not a theorem. \status{C} conjecture.

\section{Setting}\label{sec:setting}

$E$ is a Euclidean space of dimension $N$ with norm $|\cdot|$; for matrices $E=\R^{K\times d}$ with the Frobenius inner
product, $N=Kd$. A norm $\nu$ on $E$ is fixed (the dual of the norm in which steps are measured), with a
\emph{direction map}: a Borel map $o:E\to E$ with $o(x)\in\partial\nu(x)$, $o(tx)=o(x)$ for $t>0$, $o(-x)=-o(x)$. The
examples are the nuclear norm with $o=\polar$ (Muon-type, orthogonalized steps; $\polar(x)=UV^\top$ for a thin singular
value decomposition $x=U\operatorname{diag}(s)V^\top$, the partial isometry on the row space when $x$ is rank-deficient),
the $\ell_1$ norm with $o=\operatorname{sign}$ (sign descent) and the Euclidean norm with $o(x)=x/|x|$ (normalized steps).
For a random $\xi\in E$ with $\E\nu(\xi)<\infty$,
\[
\psi_\xi(G)=\E[\nu(G+\xi)-\nu(\xi)],\qquad \Med_\nu(\xi)=\arg\min_{m}\E[\nu(\xi-m)-\nu(\xi)]=-\arg\min\psi_\xi .
\]
$\psi_\xi$ is convex and $\psi_\xi(G)\ge\nu(G)-2\E\nu(\xi)$, so $\arg\min\psi_\xi$ is nonempty, convex and compact, and
$\E\,o(G+\xi)\in\partial\psi_\xi(G)$. For $\Sigma\succ0$ and $Z_\Sigma\sim N(0,\Sigma)$ put $\psi_\Sigma=\psi_{Z_\Sigma}$,
$\Phi_\Sigma=\nabla\psi_\Sigma$, and
\[
J_\Sigma=D\Phi_\Sigma(0)=\nabla^2\psi_\Sigma(0),\qquad T_\Sigma[K_3]=D^3\Phi_\Sigma(0)[K_3]=\nabla^4\psi_\Sigma(0)[K_3,\cdot\,],\qquad
\gs=\gs(\Sigma,K_3)=-\tfrac16J_\Sigma^{-1}T_\Sigma[K_3]
\]
for a symmetric $3$-tensor $K_3$. Throughout, $\gs$ is the \emph{bias coefficient}: the claim of the first-order law is
$\arg\min\psi_{\xi_b}=\gs/b+o(1/b)$ for the mean $\xi_b$ of $b$ samples with covariance $\Sigma$ and third cumulant $K_3$.

% Module: Gaussian side (used by both papers). Labels: lem:gauss
\begin{lemma}[Gaussian side]\label{lem:gauss}
Let $\nu$ be a norm on $E$ with direction map $o$, and $\Sigma\succ0$. Write $\varphi_\Sigma$ for the density of
$Z_\Sigma\sim N(0,\Sigma)$ and $\He^\Sigma_k=(-1)^k\varphi_\Sigma^{-1}D^k\varphi_\Sigma$ for the Hermite tensors.
\begin{enumerate}
\item $\psi_\Sigma$ is $C^\infty$, $\Phi_\Sigma=\nabla\psi_\Sigma$ is odd, and the derivatives of $\Phi_\Sigma$ of every
order are bounded on $E$, uniformly for $\Sigma$ in a compact set of positive definite matrices.
\item For every symmetric $k$-tensor $\cK$ and every $x\in E$,
$\int o(x+y)\inner{\cK}{\He^\Sigma_k(y)}\varphi_\Sigma(y)\,dy=D^k\Phi_\Sigma(x)[\cK]$. In particular
$J_\Sigma=\E[o(Z_\Sigma)\otimes\Sigma^{-1}Z_\Sigma]$ and $T_\Sigma[\cK]=\E[o(Z_\Sigma)\inner{\cK}{\He_3^\Sigma(Z_\Sigma)}]$.
\item The even-order derivatives of $\Phi_\Sigma$ vanish at $0$.
\item $J_\Sigma$ is symmetric positive definite, and $\lambda_{\min}(J_\Sigma)$ is bounded below on compact sets of
positive definite $\Sigma$.
\item $J_{c\Sigma}=c^{-1/2}J_\Sigma$ and $T_{c\Sigma}=c^{-3/2}T_\Sigma$ for $c>0$; hence $\gs(c\Sigma,cK_3)=\gs(\Sigma,K_3)$
and $\gs$ is linear in $K_3$.
\item There is $c_0>0$, uniform on compact sets of $\Sigma$, with $|\Phi_\Sigma(h)|\ge c_0\min(|h|,1)$.
\end{enumerate}
\end{lemma}
\begin{proof}
A norm is Lipschitz, hence differentiable almost everywhere, with gradient $o$ wherever it is differentiable; so
$\nabla\psi_\Sigma(h)=\E\,o(h+Z_\Sigma)=\int o(y)\varphi_\Sigma(y-h)\,dy$ by dominated convergence, and one may
differentiate under the integral any number of times because $o$ is bounded: the $k$-th derivative is bounded by
$\sup|o|\int|D^k\varphi_\Sigma|$. This is (1), and $D_h^k[\varphi_\Sigma(y-h)]=(-1)^kD^k\varphi_\Sigma(y-h)
=\He_k^\Sigma(y-h)\varphi_\Sigma(y-h)$ gives (2) after the change of variable $y\mapsto x+y$. Oddness of $o$ and the
symmetry of $Z_\Sigma$ make $\Phi_\Sigma$ odd, whence (3). (5) follows from $o(tx)=o(x)$:
$\Phi_{c\Sigma}(h)=\Phi_\Sigma(h/\sqrt c)$.

(4) $J_\Sigma=\nabla^2\psi_\Sigma(0)$ is symmetric. Fix $v\ne0$, let
$s=\inner{\Sigma^{-1}v}{Z}/\inner v{\Sigma^{-1}v}$ and $Z'=Z-sv$; these are independent and $s\sim N(0,\tau^2)$ with
$\tau^{-2}=\inner v{\Sigma^{-1}v}$. Then $\psi_\Sigma(tv)+\E\nu(Z)=\E\,\hbar(s+t)$ with $\hbar(\varrho)=\E\nu(Z'+\varrho v)$
convex and $|\varrho|\nu(v)-\E\nu(Z')\le\hbar(\varrho)\le\E\nu(Z')+|\varrho|\nu(v)$. Hence the nondecreasing function
$\hbar'$ goes from $-\nu(v)$ to $\nu(v)$, the measure $d\hbar'$ has total mass $2\nu(v)>0$, and
$\inner{J_\Sigma v}v=\int\varphi_{\tau}\,d\hbar'>0$ because the Gaussian density $\varphi_\tau$ is positive. The lower
bound is uniform by continuity of $\Sigma\mapsto J_\Sigma$ and compactness.

(6) For a unit vector $\hat h$, $t\mapsto\inner{\Phi_\Sigma(t\hat h)}{\hat h}$ is nondecreasing by convexity of
$\psi_\Sigma$, and equals $t\inner{J_\Sigma\hat h}{\hat h}+O(t^3)$ near $0$ by (1) and (3), with uniform constants.
\end{proof}


\section{Map of the results}\label{sec:deps}

{\small
\begin{longtable}{@{}p{0.16\textwidth}p{0.36\textwidth}p{0.24\textwidth}p{0.16\textwidth}@{}}
\toprule
Result & Content & Depends on & Status / source\\
\midrule
\endhead
Lemma~\ref{lem:gauss} & Gaussian side: smoothness, $J_\Sigma\succ0$, parity, Hermite identity & -- & \status{P} note 1\\
Thm.~\ref{thm:conv} & momentum iteration converges a.s.\ to $\arg\min\psi_{N_\beta}$; any norm, any $L^2$ noise & -- & \status{P} note 3\\
Thm.~\ref{thm:cf} & first-order law under Cram\'er's condition & Lem.~\ref{lem:gauss}, \ref{lem:osc}, \ref{lem:edge}; \citet[Thm.~20.1]{br} & \status{P*} note 1\\
Thm.~\ref{thm:beff} & linear filters: $\rho_3/\rho_2$; skewness-effective batch & Thm.~\ref{thm:cf} & \status{P*} notes 1, 2\\
Thm.~\ref{thm:state} & parameter-dependent noise: zeros at $W_\ast+H^{-1}\gs/b$ & Lem.~\ref{lem:gauss}, uniform expansion & \status{P*} note 1\\
Thm.~\ref{thm:two} & two-moment theorem for the nuclear norm & Lem.~\ref{lem:count}; divergence formula \citep{candes} & \status{P*} note 1\\
Thm.~\ref{thm:a1} & $a_1<0$, i.e.\ $a_1+a_2>3/(K+d+1)$ & Thm.~\ref{thm:two}(b); exact arithmetic & \status{CA} note 2\\
Cor.~\ref{cor:dim}--\ref{cor:compare} & dimension factor, rank-one noise, no domination & Thm.~\ref{thm:two}, \ref{thm:a1} & \status{P}\\
\midrule
Def.~\ref{def:class} & class $\mathcal X(\varepsilon)$ (cube expansion) & -- & new formulation\\
Lem.~\ref{lem:test}, \ref{lem:ball} & test functions; small balls & (E1)--(E3) & \status{P} note 2\\
Prop.~\ref{prop:sums} & weighted bounded sums are in the class, no Cram\'er & Lem.~\ref{lem:cube}, \ref{lem:near} & \status{P} note 2\\
Prop.~\ref{prop:worclass} & samples without replacement are in the class & \citet{hoeffding} & \status{P*} note 3\\
Thm.~\ref{thm:codim} & codimension $\ge2$: law holds, sharp rates & Prop.~\ref{prop:field}, Lem.~\ref{lem:locate} & \status{P} note 2\\
Thm.~\ref{thm:laplace} & Fourier decay of $\Delta\nu$: law holds; all subgradients & Prop.~\ref{prop:value}, Lem.~\ref{lem:upgrade} & \status{P} note 3\\
Lem.~\ref{lem:matrices}, \ref{lem:twobytwo} & decay for invariant norms ($k,l\ge3$) and for $2\times2$ & van der Corput \citep{stein} & \status{P*} notes 3, here\\
Thm.~\ref{thm:nuclear} & nuclear norm, every size but $1\times1$ & Thm.~\ref{thm:codim}, \ref{thm:laplace} & \status{P*}\\
Thm.~\ref{thm:sign} & sign descent: lattice sawtooth & Esseen \citep{esseen,gk} & \status{P*} note 2\\
Thm.~\ref{thm:wor} & without replacement: factor $\frac{n-2b}{n-2}$ & Prop.~\ref{prop:worclass} & \status{P*} note 3\\
Cor.~\ref{cor:mom} & batch and momentum, uniformly in $\beta$ & Thm.~\ref{thm:codim}, \ref{thm:laplace} & \status{P*}\\
Prop.~\ref{prop:lattice2} & lattice term of order $b^{-3/2}$ & -- & \status{F} note 2\\
\bottomrule
\end{longtable}}

\paragraph{Logical structure.} There are three layers. \emph{(1) Dynamics:} Theorem~\ref{thm:conv} reduces the question
``where does the iteration go'' to the location of $\arg\min\psi_{N_\beta}$, a statement about a law. \emph{(2) Location:}
$\arg\min\psi_{\xi}\approx(\rho_3/\rho_2)\gs$ for averages $\xi$ of many samples; this is proved under Cram\'er's
condition (Theorem~\ref{thm:cf}) and, with rates and no arithmetic condition, for bounded noise by two independent
mechanisms (Theorems~\ref{thm:codim} and~\ref{thm:laplace}), which between them cover the nuclear norm in every size
except $1\times1$, where the law fails (Theorem~\ref{thm:sign}). \emph{(3) Evaluation:} for the nuclear norm and isotropic
covariance $\gs$ is explicit (Theorems~\ref{thm:two} and~\ref{thm:a1}). Layers (1) and (3) are independent of each other
and of the mechanism used in (2).

\section{Dynamics: the iteration converges to a norm median}\label{sec:dyn}

\begin{theorem}[convergence with momentum and arbitrary noise]\label{thm:conv}\StmtConv\end{theorem}
% Module: convergence of norm-steepest descent with momentum, arbitrary square-integrable noise.
% Label: proof of thm:conv
\begin{proof}[Proof of Theorem~\ref{thm:conv}]
Let $K_\nu$ be the Lipschitz constant of $\nu$ for the Euclidean norm, so that $|O_t|\le K_\nu$, and let $L$, $\mu$ be
the smoothness and strong convexity constants of $F$. Extend the innovations to an i.i.d.\ sequence $(\xi_t)_{t\in\Z}$
and put $N_t=(1-\beta)\sum_{j\ge0}\beta^j\xi_{t-j}$: a stationary, square-integrable Markov chain with
$N_t=\beta N_{t-1}+(1-\beta)\xi_t$ and marginal law $\pi$, the law of $N_\beta$. Let $\mathcal F_t=\sigma(\xi_s,s\le t)$.
By induction $W_t$ and $G_t=\nabla F(W_t)$ are $\mathcal F_{t-1}$-measurable.

\emph{Step 1: the buffer tracks $G_t+N_t$.} Let $e_t=M_t-G_t-N_t$. Then $e_1=\beta(M_0-G_1-N_0)$ and, for $t\ge2$,
$e_t=\beta e_{t-1}+\beta(G_{t-1}-G_t)$ with $|G_t-G_{t-1}|\le LK_\nu\eta_{t-1}$. Hence
$|e_t|\le\beta^{t-1}|e_1|+LK_\nu\sum_{s<t}\beta^{t-s}\eta_s$ and, the lengths being nonincreasing,
\begin{equation}\label{eq:esum}
\sum_t\eta_t|e_t|\le|e_1|\,\eta_1\sum_t\beta^{t-1}+LK_\nu\sum_s\eta_s^2\sum_{t>s}\beta^{t-s}<\infty\qquad\text{almost surely.}
\end{equation}

\emph{Step 2: a pathwise inequality on loss differences.} Fix $a\in\mathcal A:=\arg\min\psi_{N_\beta}$, a nonempty convex
compact set, and let
\[
f_g(n)=\nu(g+n)-\nu(a+n),\qquad\hbar(g)=\int f_g\,d\pi=\psi_{N_\beta}(g)-\psi_{N_\beta}(a)\ \ge0 .
\]
The subgradient inequality for $\nu$ at the \emph{actual} argument $M_t=G_t+N_t+e_t$, tested against $a+N_t$, reads
$\nu(a+N_t)\ge\nu(M_t)+\inner{O_t}{a-G_t-e_t}$, whence
\begin{equation}\label{eq:path}
\inner{O_t}{G_t-a}\ \ge\ \nu(M_t)-\nu(a+N_t)-K_\nu|e_t|\ \ge\ f_{G_t}(N_t)-2K_\nu|e_t| .
\end{equation}
No regularity of the law of the noise and no choice of the subgradient enter here.

\emph{Step 3: Poisson equation.} Let $Q$ be the transition operator of $(N_t)$. For $j\ge0$,
$Q^jf_g(n)=\E f_g(\beta^jn+N^{(j)})$ with $N^{(j)}=(1-\beta)\sum_{i<j}\beta^i\xi_i$, and
$\hbar(g)=\E f_g(\beta^jN'+N^{(j)})$ with $N'\sim\pi$ independent. The function $n\mapsto f_g(n)$ is $2K_\nu$-Lipschitz
and $g\mapsto f_g(n)$ is $K_\nu$-Lipschitz, uniformly. Therefore, with $m_1=\E|N_\beta|$,
\[
\big|Q^jf_g(n)-\hbar(g)\big|\le2K_\nu\beta^j(|n|+m_1),\qquad
\big|[Q^jf_g(n)-\hbar(g)]-[Q^jf_{g'}(n)-\hbar(g')]\big|\le2K_\nu|g-g'| .
\]
Define $U_g(n)=\sum_{j\ge0}(Q^jf_g(n)-\hbar(g))$ and $V_g=QU_g=\sum_{j\ge1}(Q^jf_g-\hbar(g))$. The series converge
absolutely, $U_g-V_g=f_g-\hbar(g)$, both are jointly continuous in $(g,n)$, and, using
$\min(x,y)\le x^\alpha y^{1-\alpha}$ termwise, for every $\alpha\in(0,1)$,
\begin{equation}\label{eq:UV}
|U_g(n)|+|V_g(n)|\le C(1+|n|),\qquad|V_g(n)-V_{g'}(n)|\le C_\alpha|g-g'|^\alpha(1+|n|)^{1-\alpha},
\end{equation}
with constants depending on $K_\nu,\beta,m_1,\alpha$ but not on $g,g'$.

\emph{Step 4: the noise terms are summable.} We claim that $\sum_t\eta_t\big(f_{G_t}(N_t)-\hbar(G_t)\big)$ converges almost
surely. Write the summand, before multiplication by $\eta_t$, as
\[
\big[U_{G_t}(N_t)-V_{G_t}(N_{t-1})\big]+\big[V_{G_t}(N_{t-1})-V_{G_t}(N_t)\big].
\]
Since $G_t$ is $\mathcal F_{t-1}$-measurable and $\xi_t$ is independent of $\mathcal F_{t-1}$,
$\E[U_{G_t}(N_t)\mid\mathcal F_{t-1}]=V_{G_t}(N_{t-1})$: the first brackets are martingale differences, with second
moments bounded by $C\,\E(1+|N_\beta|)^2$ uniformly in $t$ by \eqref{eq:UV} (no bound on $G_t$ is needed). As
$\sum_t\eta_t^2<\infty$, their weighted sum converges almost surely. For the second brackets, summation by parts gives
\[
\sum_{t=1}^T\eta_t\big[V_{G_t}(N_{t-1})-V_{G_t}(N_t)\big]=\eta_1V_{G_1}(N_0)-\eta_TV_{G_T}(N_T)
+\sum_{t=1}^{T-1}\big[\eta_{t+1}V_{G_{t+1}}(N_t)-\eta_tV_{G_t}(N_t)\big].
\]
The boundary term $\eta_TV_{G_T}(N_T)$ tends to $0$ almost surely because $\sum_T\eta_T^2\E(1+|N_T|)^2<\infty$. By
\eqref{eq:UV} and $|G_{t+1}-G_t|\le LK_\nu\eta_t$, the $t$-th internal difference is at most
$C(\eta_t-\eta_{t+1})(1+|N_t|)+C_\alpha\eta_t^{1+\alpha}(1+|N_t|)^{1-\alpha}$ in absolute value; choosing
$\alpha\in(\frac1\theta-1,1)$, which is possible because $\theta>\frac12$, the expectations are summable. This proves the
claim.

\emph{Step 5: Lyapunov function.} Let $W^a=(\nabla F)^{-1}(a)$ (strong convexity makes $\nabla F$ a bijection of $E$ with
a $\frac1\mu$-Lipschitz inverse) and $B_t(a)=F(W_t)-F(W^a)-\inner a{W_t-W^a}\ge\frac\mu2|W_t-W^a|^2$. Smoothness of $F$ and
\eqref{eq:path} give
\[
B_{t+1}(a)\le B_t(a)-\eta_t\inner{O_t}{G_t-a}+\tfrac L2K_\nu^2\eta_t^2
\le B_t(a)-\eta_t\hbar(G_t)-\eta_t\big(f_{G_t}(N_t)-\hbar(G_t)\big)+2K_\nu\eta_t|e_t|+\tfrac L2K_\nu^2\eta_t^2 .
\]
By \eqref{eq:esum} and Step 4 the series of the last three terms converges almost surely; subtracting its partial sums
from $B_t(a)+\sum_{s<t}\eta_s\hbar(G_s)$ gives a nonincreasing sequence bounded below. Hence $B_t(a)$ converges and
$\sum_t\eta_t\hbar(G_t)<\infty$. In particular $(W_t)$ and $(G_t)$ are bounded and, as $\sum_t\eta_t=\infty$, there is a
cluster point $\widehat G$ of $(G_t)$ with $\hbar(\widehat G)=0$, i.e.\ $\widehat G\in\mathcal A$.

\emph{Step 6: a single limit.} Almost surely the above holds simultaneously for all $a$ in a countable dense subset of
$\mathcal A$. On bounded sets of $W$ the maps $a\mapsto F(W)-F(W^a)-\inner a{W-W^a}$, $a\in\mathcal A$, are equicontinuous,
so $B_t(a)$ converges for every $a\in\mathcal A$, in particular for $a=\widehat G$. Along the subsequence defining
$\widehat G$, $W_t\to W^{\widehat G}$ and $B_t(\widehat G)\to0$; hence $B_t(\widehat G)\to0$ along the whole sequence, and
$W_t\to W^{\widehat G}$, $G_t\to\widehat G$.
\end{proof}


\begin{remark}
For $\beta=0$ and the nuclear norm this is Theorem~2(b) of the original manuscript (stochastic polar steps are
stochastic mirror descent in the space of gradients), whose proof is the special case $e_t=0$, $U_g=f_g-\hbar(g)$. The
first note proved the momentum case under a small-ball assumption on the noise, which excludes discrete noise;
the present proof replaces the comparison of polar factors by a comparison of loss values, which is insensitive to
atoms. The limit is random when $\arg\min\psi_{N_\beta}$ is not a point (example: $\nu=|\cdot|$ on $\R$, $\xi=\pm1$,
$\beta=\frac13$, where the set is $[-\frac13,\frac13]$).
\end{remark}

\section{Location under Cram\'er's condition}\label{sec:cramer}

\begin{assumption}[smooth noise]\label{ass:N}\StmtAssN\end{assumption}
\begin{theorem}[first-order law]\label{thm:cf}\StmtCF\end{theorem}
% Module: first-order law under Cramer's condition. Labels: lem:osc, lem:edge, proof of thm:cf, proof of thm:beff
\begin{lemma}[oscillation]\label{lem:osc}
For the $\ell_1$, Euclidean and nuclear norms with their direction maps, and every compact set $\mathcal C$ of positive
definite covariances, there is $C$ such that for $\varepsilon\le1$
\[
\sup_{\Sigma\in\mathcal C}\ \sup_{y\in E}\ \E\,\mathrm{osc}\big(o;B(y+Z_\Sigma,\varepsilon)\big)\le C\varepsilon\log(e/\varepsilon),
\]
where $\mathrm{osc}(o;B)$ is the largest oscillation on $B$ of a coordinate of $o$. Moreover
$\sup_y\Prob(\sigma_{\min(K,d)}(y+Z_\Sigma)\le t)\le Ct$.
\end{lemma}
\begin{proof}
$\ell_1$: the oscillation vanishes unless some coordinate of $y+Z$ is within $\varepsilon$ of $0$, an event of
probability at most $N\varepsilon\sup\varphi$. Euclidean: $\mathrm{osc}\le\min(2,2\varepsilon/(|x|-\varepsilon)_+)$.
Nuclear, $K\le d$ (transpose otherwise): with rows $\varrho_j$ of $X$,
$\sigma_K(X)\ge K^{-1/2}\min_j\dist(\varrho_j,\operatorname{span}\{\varrho_l\}_{l\ne j})$; conditionally on the other rows,
$\varrho_j$ is Gaussian with covariance at least $\lambda_{\min}(\Sigma)I$ and its component along a unit normal to the
span has a density bounded by $(2\pi\lambda_{\min})^{-1/2}$; a union bound gives the small-ball estimate. By
\citet{lisun}, $\nF{\polar(A)-\polar(B)}\le2\nF{A-B}/(\sigma_K(A)+\sigma_K(B))$ for $A,B$ of rank $K$; hence
$\mathrm{osc}(\polar;B(x,\varepsilon))\le\min(2,4\varepsilon/\sigma_K(x))$ when $\sigma_K(x)>2\varepsilon$, and
$\E\min(2,4\varepsilon/\sigma_K)\le C\varepsilon\log(e/\varepsilon)$.
\end{proof}

\begin{lemma}[one-term expansion of the mean field, uniform in the shift]\label{lem:edge}
Under Assumption~\ref{ass:N}, for the norms of Lemma~\ref{lem:osc}, with $X_b=\sqrt b\,\xi_b$,
\[
\sup_{h\in E}\Big|\E\,o(X_b+h)-\Phi_\Sigma(h)-\frac1{6\sqrt b}\,D^3\Phi_\Sigma(h)[K_3]\Big|=o(b^{-1/2}).
\]
\end{lemma}
\begin{proof}
Theorem~20.1 of \citet{br} with $s=3$ provides $\delta_b=o(b^{-1/2})$ and $c,c'>0$, depending on the law of $\zeta$ only,
such that for every bounded Borel $f$,
$|\int f\,d(Q_b-\Phi_{0,\Sigma}-b^{-1/2}P_1)|\le\sup|f|\,\delta_b+c'\sup_y\int\mathrm{osc}(f;B(x+y,2e^{-cb}))\varphi_\Sigma(x)dx$,
where $Q_b$ is the law of $X_b$ and $P_1$ has density $\frac16\inner{K_3}{\He_3^\Sigma}\varphi_\Sigma$. Apply it to the
coordinates of $f=o(\cdot+h)$: the supremum is bounded, the oscillation term is $O(be^{-cb})$ by Lemma~\ref{lem:osc},
uniformly in $h$, and $\int o(x+h)\frac16\inner{K_3}{\He^\Sigma_3(x)}\varphi_\Sigma(x)dx=\frac16D^3\Phi_\Sigma(h)[K_3]$ by
Lemma~\ref{lem:gauss}(2).
\end{proof}

\begin{proof}[Proof of Theorem~\ref{thm:cf}]
By homogeneity $\Phi_{\xi_b}(G)=\E\,o(X_b+\sqrt b\,G)$. Let $G=(\gs+u)/b$ with $|u|\le R$, so $h=\sqrt bG=O(b^{-1/2})$.
By Lemma~\ref{lem:gauss}, $\Phi_\Sigma(h)=J_\Sigma h+O(|h|^3)$ and $D^3\Phi_\Sigma(h)[K_3]=T_\Sigma[K_3]+O(|h|^2)$; with
Lemma~\ref{lem:edge},
\[
\Phi_{\xi_b}(G)=b^{-1/2}\big(J_\Sigma(\gs+u)+\tfrac16T_\Sigma[K_3]\big)+o(b^{-1/2})=b^{-1/2}J_\Sigma u+o(b^{-1/2}),
\]
uniformly in $|u|\le R$. Take $|u|=\epsilon$: $\inner{\Phi_{\xi_b}(G)}{u}\ge b^{-1/2}(\lambda_{\min}(J_\Sigma)\epsilon^2-o(1))>0$
for $b\ge b_0(\epsilon)$. If a minimizer $G'$ of $\psi_{\xi_b}$ had $|G'-\gs/b|>\epsilon/b$, let $G$ be the point of the
segment $[\gs/b,G']$ on the sphere of radius $\epsilon/b$, so that $G'-G=\tau u$ with $\tau>0$; since
$\Phi_{\xi_b}(G)\in\partial\psi_{\xi_b}(G)$, $\psi_{\xi_b}(G')\ge\psi_{\xi_b}(G)+\tau\inner{\Phi_{\xi_b}(G)}u>\psi_{\xi_b}(G)$,
a contradiction.
\end{proof}

\begin{proof}[Proof of Theorem~\ref{thm:beff}]
Write the filtered noise as a batch mean of regrouped samples:
$N_w=\sum_kw_k\xi_b^{(k)}=\frac1b\sum_{i\le b}Y_i$ with $Y_i=\sum_kw_k\zeta_{k,i}$ i.i.d. The series converges absolutely in
$L^3$ because $\sum_k|w_k|<\infty$; cumulants add over independent summands, so $Y$ has mean $0$, covariance
$\rho_2\Sigma$ and third cumulant $\rho_3K_3$; and $|\E e^{i\inner tY}|=\prod_k|\varphi(w_kt)|\le|\varphi(w_{k_0}t)|$ for an
index with $w_{k_0}\ne0$, so $Y$ satisfies Cram\'er's condition. Theorem~\ref{thm:cf} applies to $Y$ and gives
$\arg\min\psi_{N_w}\subset B\big(\gs(\rho_2\Sigma,\rho_3K_3)/b,\,o(1/b)\big)$, and
$\gs(\rho_2\Sigma,\rho_3K_3)=\frac{\rho_3}{\rho_2}\gs(\Sigma,K_3)$ by Lemma~\ref{lem:gauss}(5). For the exponential
average, $\rho_p=(1-\beta)^p/(1-\beta^p)$ and $\rho_3/\rho_2=(1-\beta^2)/(1+\beta+\beta^2)$.
\end{proof}

\begin{theorem}[linear filters and the skewness-effective batch size]\label{thm:beff}\StmtBeff\end{theorem}
(The proof is at the end of the preceding block. It replaces the Lindeberg argument of the first note; three moments
suffice.)

\begin{theorem}[parameter-dependent noise]\label{thm:state}\StmtState\end{theorem}
\begin{proposition}[linear noise models]\label{prop:lin}\StmtLin\end{proposition}
% Module: parameter-dependent noise. Label: eq:hb
\begin{proof}[Proof of Theorem~\ref{thm:state}]
By homogeneity and (iii), uniformly in $W\in\mathcal U$,
\begin{equation}\label{eq:hb}
h_b(W)=\Phi_{\Sigma(W)}\big(\sqrt b\,\nabla F(W)\big)+\tfrac1{6\sqrt b}D^3\Phi_{\Sigma(W)}\big(\sqrt b\,\nabla F(W)\big)[K_3(W)]+o(b^{-1/2}).
\end{equation}
\emph{Step 1.} By Lemma~\ref{lem:gauss}, uniformly over the compact set of covariances,
$|\Phi_{\Sigma(W)}(h)|\ge c_0\min(|h|,1)$ and the second term of \eqref{eq:hb} is at most $C_1b^{-1/2}$. Hence
$|h_b(W)|\le\delta b^{-1/2}$ forces $c_0\min(\sqrt b|\nabla F(W)|,1)\le(C_1+\delta+o(1))b^{-1/2}$, i.e.\
$|\nabla F(W)|\le C_2/b$ and $|W-W_\ast|\le C_3/b$ for $b$ large. \emph{Step 2.} Write $W=W_\ast+w/b$, $|w|\le C_3$. Then
$\sqrt b\,\nabla F(W)=b^{-1/2}(Hw+o(1))$, $\Sigma(W)\to\Sigma_\ast$, $K_3(W)\to K_{3\ast}$, and $J_\Sigma$, $T_\Sigma$ are
continuous in $\Sigma$, so \eqref{eq:hb} reads $h_b(W)=b^{-1/2}\big(J_{\Sigma_\ast}Hw+\frac16T_{\Sigma_\ast}[K_{3\ast}]+o(1)\big)$,
uniformly in $|w|\le C_3$. If $|h_b(W)|\le\delta b^{-1/2}$ then $|w-H^{-1}\gs|\le\|(J_{\Sigma_\ast}H)^{-1}\|(\delta+o(1))$.
\emph{Existence.} The continuous field $w\mapsto\sqrt b\,h_b(W_\ast+w/b)$ converges uniformly on a ball around
$H^{-1}\gs$ to the affine isomorphism $w\mapsto J_{\Sigma_\ast}H(w-H^{-1}\gs)$, which has degree $\pm1$ there.
\end{proof}

\begin{proof}[Proof of Proposition~\ref{prop:lin}]
Apply Theorem~20.1 of \citet{br} to $\chi$ and to $f_{W,h}(x)=o(L_Wx+h)$: its oscillation on a ball of radius
$\varepsilon$ is that of $o$ on a ball of radius $\|L_W\|\varepsilon$ around $L_Wx+h$, and $L_WZ_\chi\sim
N(0,\Sigma(W))$ with $\Sigma(W)=L_W\Sigma_\chi L_W^\top$ in a compact set of positive definite matrices, so
Lemma~\ref{lem:osc} bounds the oscillation term uniformly; the remainder $\delta_b$ depends on $\chi$ only. The
Gaussian and first Edgeworth terms of $\chi$ push forward to those of covariance $\Sigma(W)$ and cumulant
$L_W^{\otimes3}K_3^\chi=K_3(W)$: with $u(s)=\E f_{W,h}(Z_\chi+s)=\Phi_{\Sigma(W)}(h+L_Ws)$,
$-\int f_{W,h}K_3^\chi\cdot\partial^3\varphi_{\Sigma_\chi}=D^3u(0)[K_3^\chi]=D^3\Phi_{\Sigma(W)}(h)[L_W^{\otimes3}K_3^\chi]$.
\end{proof}


\section{Evaluation: the two-moment theorem}\label{sec:two}

\begin{theorem}[two-moment theorem]\label{thm:two}\StmtTwo\end{theorem}
% Module: two-moment theorem (proof), counting lemma, proofs of corollaries. Labels: lem:count
\begin{proof}[Proof of Theorem~\ref{thm:two}]
\emph{(a) Reduction.} $J_I$ and $T_I$ are equivariant under $Z\mapsto UZV^\top$, $U\in O(K)$, $V\in O(d)$, because
$\polar(UZV^\top)=U\polar(Z)V^\top$ and the law of $Z$ is invariant. $E=\R^K\otimes\R^d$ is irreducible, so
$J_I=c\,\mathrm{Id}$ with $c=\tr J_I/N=\E\inner{\polar(Z)}Z/N=c_{K,d}$ \citep{xie2026muon,chen2026when}. $T_I$ is an invariant element of
$E^{\otimes4}=(\R^K)^{\otimes4}\otimes(\R^d)^{\otimes4}$; the invariants of $O(K)\times O(d)$ are spanned by the
products $\pi\otimes\pi'$ of an $O(K)$-invariant and an $O(d)$-invariant, and by the first fundamental theorem for the
orthogonal group each factor is spanned by the three pairings of the four indices. Label a pairing by the input slot
paired with the output index; $T_I$ is symmetric in its three input slots, the symmetric group $S_3$ permutes the nine
products $(\pi,\pi')$ with two orbits, $\{\pi=\pi'\}$ and $\{\pi\ne\pi'\}$, and the orbit averages are the tensors of
$\cA$ and $\cB$ on symmetric arguments. Hence $T_I=\alpha\cA+\beta\cB$. To find $\alpha,\beta$ contract with one representative of each orbit:
using $\He_3(Z)^{(ab),(cd),(ef)}=Z_{ab}Z_{cd}Z_{ef}-Z_{ab}\delta_{(cd),(ef)}-Z_{cd}\delta_{(ab),(ef)}-Z_{ef}\delta_{(ab),(cd)}$,
\begin{align*}
\textstyle\sum_{ij}\cA\big[\E\,\polar(Z)_{ij}\He_3(Z)\big]_{ij}&=\E\big[\nnuc Z(\nF Z^2-N-2)\big]=-m_1,\\
\textstyle\sum_{ij}\cB\big[\E\,\polar(Z)_{ij}\He_3(Z)\big]_{ij}&=\E\big[\inner{\polar(Z)}{ZZ^\top Z}\big]-(K+d+1)m_1=m_3-(K+d+1)m_1,
\end{align*}
where $\E[\nnuc Z\nF Z^2]=(N+1)m_1$ because $\nF Z$ is independent of $Z/\nF Z$ and $\E\chi_N^3=(N+1)\E\chi_N$. On the
other side, by Lemma~\ref{lem:count}, the same contractions of $\alpha\cA+\beta\cB$ are
$\frac N3[(N+2)\alpha+(K+d+1)\beta]$ and $\frac N3[(K+d+1)\alpha+\frac12(N+K+d+3)\beta]$. The determinant of this system
is $\frac{N^2}{18}(K-1)(d-1)(K+2)(d+2)$; solving, and substituting $m_3=(K+d-q)m_1$ from (b), gives
$\alpha=-c_{K,d}a_1$, $\beta=-c_{K,d}a_2$ with \eqref{eq:a12}. Then \eqref{eq:two} is Theorem~\ref{thm:cf} with the scaling
of Lemma~\ref{lem:gauss}. The map $K_3\mapsto(\cA[K_3],\cB[K_3])$ is onto $E\oplus E$ (the two equivariant maps are
linearly independent and $E$ is irreducible), so $K_3\mapsto a_1\cA+a_2\cB$ is onto $E$.

\emph{(b)} $\sum_is_i^3=\inner Z{f(Z)}$ with the spectral function $f(Z)=(ZZ^\top)^{1/2}Z=U\operatorname{diag}(s_i^2)V^\top$,
which is $C^1$; by Stein's identity $m_3=\E\operatorname{div}f(Z)$ and by the divergence formula for spectral functions
\citep{candes}, with $\varphi(s)=s^2$,
$\operatorname{div}f=\sum_i\big(\varphi'(s_i)+|K-d|\frac{\varphi(s_i)}{s_i}\big)+\sum_{i\ne j}\frac{s_i\varphi(s_i)-s_j\varphi(s_j)}{s_i^2-s_j^2}
=(K+d)\sum_is_i-\sum_{i\ne j}\frac{s_is_j}{s_i+s_j}$, since
$\frac{s_i^3-s_j^3}{s_i^2-s_j^2}=s_i+s_j-\frac{s_is_j}{s_i+s_j}$. The bounds on $q$ follow from
$0<\frac{s_is_j}{s_i+s_j}<\frac{s_i+s_j}4$ almost surely, and $a_2>0$ from \eqref{eq:a12}.

\emph{(c)} Algebra from \eqref{eq:a12}, using $Kd+1-(K+d)=(K-1)(d-1)$ and, for $K\le d$,
$Kd+2-K(K+d+1)=-(K-1)(K+2)$; positivity from $q<\frac{r-1}2\le\frac{\bar d-1}2$.

\emph{(d)} $\nabla^4\psi_I(0)[D,D,D,D]=\inner D{T_I[D^{\otimes3}]}$, $\cA[D^{\otimes3}]=\nF D^2D$,
$\cB[D^{\otimes3}]=DD^\top D$. If $a_1\ge0$ the sign is clear. If $a_1<0$, $\nF D^4\le\operatorname{rank}(D)\tr(DD^\top)^2$
gives $a_1\nF D^4+a_2\tr(DD^\top)^2\ge(r\,a_1+a_2)\tr(DD^\top)^2>0$ by (c).

\emph{(e)} The empirical distribution of the $s_i^2/d$ converges almost surely to the Marchenko--Pastur law of ratio
$\gamma$, with convergence of the moments $m_1/(K\sqrt d)$ and $m_3/(Kd^{3/2})$ (the extreme singular values are bounded in
every $L^p$); by (b), $q/K=(K+d)/K-m_3/(Km_1)\to[(1+\gamma)\mu_{1/2}-\mu_{3/2}]/(\gamma\mu_{1/2})$. The hypergeometric
form of $\mu_s$ is the Narayana expansion for integer $s$; for $s=\frac12,\frac32$ only the integral definition of $\mu_s$ is used; $\mu_{1/2}(1)=\frac8{3\pi}$, $\mu_{3/2}(1)=\frac{64}{15\pi}$
(quarter-circle law), hence $\kappa(1)=2-\frac85$. The limits of $a_1,a_2$ follow from \eqref{eq:a12}.

\emph{$K=1$.} The same computation with a single invariant; $m_3=(d+1)m_1$ and the contraction of the invariant with
itself is $N(N+2)/3$, so $\alpha+\beta=-3m_1/(d(d+2))$.
\end{proof}

\begin{lemma}[counting]\label{lem:count}
Let $\mathsf A_1$, $\mathsf B_1$ be one pairing tensor from each orbit and $\bar{\mathsf A}$, $\bar{\mathsf B}$ the orbit
averages (so that $\bar{\mathsf A}$, $\bar{\mathsf B}$ are the tensors of $\cA$, $\cB$ on symmetric arguments). Then
$\inner{\mathsf A_1}{\bar{\mathsf A}}=\frac{N(N+2)}3$, $\inner{\mathsf A_1}{\bar{\mathsf B}}=\inner{\mathsf B_1}{\bar{\mathsf A}}=\frac{N(K+d+1)}3$,
$\inner{\mathsf B_1}{\bar{\mathsf B}}=\frac{N(N+K+d+3)}6$.
\end{lemma}
\begin{proof}
The inner product of two pairing tensors is $K^{\ell}d^{\ell'}$, $\ell,\ell'\in\{1,2\}$ being the numbers of closed
loops of row and of column indices. The three inner products of $\mathsf A_1$ with the first orbit are $N^2,N,N$; its
six inner products with the second orbit are $Kd^2,Kd^2,K^2d,K^2d,Kd,Kd$; the six inner products of $\mathsf B_1$ with
its orbit are $K^2d^2,K^2d,Kd^2,Kd,Kd,Kd$. (Direct enumeration; also verified by brute-force contraction for seven sizes.)
\end{proof}

\begin{proof}[Proof of Corollary~\ref{cor:dim}]
For independent entries the only nonzero components of $K_3$ are $K_3^{(ij),(ij),(ij)}=\kappa_3$, so
$\cA[K_3]_{ij}=\cB[K_3]_{ij}=\kappa_3$. The value of $a_1+a_2$ is Theorem~\ref{thm:two}(c); the upper bound is
$q_{K,d}<\frac{r-1}2$ with $(r+2)(\bar d+2)=(K+2)(d+2)$, and the lower bound is Theorem~\ref{thm:a1}.
\end{proof}

\begin{proof}[Proof of Corollary~\ref{cor:rank1}]
If $\zeta=uv^\top$ then $\zeta\zeta^\top\zeta=|u|^2|v|^2uv^\top=\nF\zeta^2\zeta$, so $\cB=\cA$ and
$\gs=(a_1+a_2)\cA/(6\sigma^2)$. For the Frobenius norm apply the case $K=1$ of Theorem~\ref{thm:two} on $\R^{Kd}$:
$\gs=\frac3{Kd+2}\cA/(6\sigma^2)$. The ratio is $\varrho=\frac{(3+2q)(Kd+2)}{(K+2)(d+2)}$; it exceeds
$\frac{Kd+2}{K+d+1}>1$ by Theorem~\ref{thm:a1} (as $Kd+2-(K+d+1)=(K-1)(d-1)>0$), and it is below
$\frac{(r+2)(Kd+2)}{(K+2)(d+2)}=\frac{Kd+2}{\bar d+2}<r$.
\end{proof}

\begin{proof}[Proof of Corollary~\ref{cor:compare}]
(i) Theorems~\ref{thm:cf} and~\ref{thm:two}; for the $\ell_1$ norm and diagonal $\Sigma$,
$J=\operatorname{diag}(\sqrt{2/\pi}/\sigma_k)$ and $T[K_3]_k=-\sqrt{2/\pi}\,\kappa_{3,k}/\sigma_k^3$, so
$(\gs)_k=\kappa_{3,k}/(6\sigma_k^2)$. (ii) The only nonzero entries of $K_3$ are the permutations of
$((11),(12),(12))$, whence $\cA=\cB=\tau e_1e_1^\top$. Such a law exists:
$\zeta_{11}=\varrho(\zeta_{12}^2-1)/\sqrt2+\sqrt{1-\varrho^2}\,\epsilon$ with $\epsilon$ an independent centred Gamma
variable whose skewness cancels that of the first term. (iii) For this tensor
$\cA_{11}=1+K_3^{(11),(22),(22)}$, $\cB_{11}=1$ and all other entries of $\cA,\cB$ vanish (the pattern
$(kj),(kl),(il)$ of $\cB$ is never a permutation of $((11),(22),(22))$), so
$a_1\cA+a_2\cB=(a_1+a_2+a_1K_3^{(11),(22),(22)})e_1e_1^\top=0$; a small multiple of this $K_3$ is the third cumulant of a law
with covariance $I$ and a density (the pairs of covariance and third moment of laws with a density near the Gaussian
fill a neighbourhood of $(I,0)$). (iv) Cumulants add over independent summands, and
$a_1\nF M^4+a_2\tr(MM^\top)^2>0$ by Theorem~\ref{thm:two}(d).
\end{proof}

\begin{theorem}[sign of $a_1$]\label{thm:a1}\StmtAone\end{theorem}
% Module: sign of a_1 (computer-assisted for K>=3). Label: eq:num
\begin{proof}[Proof of Theorem~\ref{thm:a1}]
Let $W=ZZ^\top$, $D_p=\E\tr W^{p/2}=\E\sum_is_i^p$, $r(p)=D_{p+2}/D_p$ and $m=K+d+1$.
\emph{Equivalences.} By the Stein identity $D_3=(K+d-q)D_1$ of Theorem~\ref{thm:two}(b), the claim is
$D_3<c^\ast D_1$ with $c^\ast=K+d-\frac{(K-1)(d-1)}{2m}$. Since $\E[\tr W\cdot\sum_js_j^p]=(Kd+p)D_p$ (differentiate the
Gaussian integral along dilations), $\E\sum_js_j^pA_j/D_p=Kd+p-r(p)$; and $r(2)=m$,
$r(4)=\E\tr W^3/\E\tr W^2=(K^2+d^2+3Kd+3K+3d+4)/m$, so $3r(2)-r(4)=2c^\ast$: this is (i) and (ii). (iii):
$\E_\varpi[x^{\pm1/2}]=D_{2\pm1}/(Kd)$, $\E_\varpi x=m$, $\E_\varpi x^2=\E\tr W^3/(Kd)$, and $m-\frac v{2m}=c^\ast$. For a Gamma law
of shape $\alpha$ and scale $\theta$ both sides of (iii) equal $\theta(\alpha-\frac12)$; for $K=1$, $\varpi$ is the $\chi^2_{d+2}$ law.

\emph{$K=2$.} The density of $(s_1>s_2)$ is proportional to $(s_1^2-s_2^2)(s_1s_2)^{d-2}e^{-(s_1^2+s_2^2)/2}$. With
$s_1+s_2=\rho\cos\frac\phi2$, $s_1-s_2=\rho\sin\frac\phi2$, $\phi\in(0,\frac\pi2)$, the variables $\rho$ and
$\gamma=\cos\phi$ are independent, $\gamma$ has density $(d-1)\gamma^{d-2}$ on $(0,1)$,
$\sum s_i=\rho\sqrt{(1+\gamma)/2}$ and $\frac{2s_1s_2}{s_1+s_2}=\rho\gamma/\sqrt{2(1+\gamma)}$. Hence
$q_{2,d}=\E[\gamma(1+\gamma)^{-1/2}]/\E[(1+\gamma)^{1/2}]$ (for $d=2$: $q_{2,2}=(3\sqrt2-2)/7$), and with $n=d-2$ the claim
is $\mathcal F=\int_0^1\gamma^n\big((n+9)\gamma-(n+1)\big)(1+\gamma)^{-1/2}d\gamma>0$. Integrating by parts,
$(n+1)\int_0^1\gamma^n(1+\gamma)^{-1/2}=2^{-1/2}+\frac12\int_0^1\gamma^{n+1}(1+\gamma)^{-3/2}$, so
$\mathcal F\ge(n+\frac{17}2)\int_0^1\gamma^{n+1}(1+\gamma)^{-1/2}d\gamma-2^{-1/2}\ge2^{-1/2}\big(\frac{n+17/2}{n+2}-1\big)>0$.

\emph{$K\ge3$.} By (iii) the claim is $\E_\varpi[g]>0$ for $g(x)=(c^\ast-x)x^{-1/2}$. For $n\ge1$, $g^{(n)}$ has the sign
$(-1)^n$ on $(0,\infty)$; so (1)~the Taylor polynomial $T_5$ of $g$ at $m$ satisfies $g\ge T_5$, and (2)~the Hermite
interpolant $H$ of $g$ at four nodes (values and first derivatives) satisfies $g\ge H$. Both $\E_\varpi T_5$ and $\E_\varpi H$
are linear in the moments $\E_\varpi x^j=\E\tr W^{j+1}/(Kd)$, $j\le7$, and $\E\tr W^k$ is a polynomial in $(K,d)$ computed by
Wick's formula (sum over the $(2k-1)!!$ pairings of $K^{\#\text{row classes}}d^{\#\text{column classes}}$). Step~1:
\begin{equation}\label{eq:num}
\begin{aligned}
\sqrt m\,\E_\varpi T_5&=\mathrm{NUM}(K,d)/(256\,m^6),\\
\mathrm{NUM}&=K^5(112d^2+112d-224)+K^4(266d^3-221d^2-3854d-3751)\\
&+K^3(266d^4-358d^3-10520d^2-33316d-36712)\\
&+K^2(112d^5-221d^4-10520d^3-52379d^2-118790d-96410)\\
&+K(112d^5-3854d^4-33316d^3-118790d^2-166352d-64872)\\
&-224d^5-3751d^4-36712d^3-96410d^2-64872d+16497,
\end{aligned}
\end{equation}
and $\mathrm{NUM}(9+\alpha,9+\alpha+\beta)$, $\mathrm{NUM}(3,22+\beta)$, $\mathrm{NUM}(4,17+\beta)$, $\mathrm{NUM}(5,14+\beta)$,
$\mathrm{NUM}(6,12+\beta)$, $\mathrm{NUM}(7,11+\beta)$, $\mathrm{NUM}(8,9+\beta)$ are polynomials in $\alpha,\beta$ with
positive coefficients. Step~2: for the $52$ remaining pairs $3\le K\le d$ ($K=3$, $d\le21$; $K=4$, $d\le16$; $K=5$,
$d\le13$; $K=6$, $d\le11$; $K=7$, $d\le10$; $(8,8)$), nodes $x_i=t_i^2$ with rational $t_i$ (so that $g(x_i),g'(x_i)$
are rational) give exact rational values $\E_\varpi H>0$, the smallest being $0.0323$ at $(3,21)$. Both steps are exact
integer computations; the moment polynomials, the identity \eqref{eq:num}, the positivity of the coefficients and the
$52$ certificates were recomputed by two independent implementations (supplementary code, \texttt{a1/} and \texttt{checks/v01, v02}).
Finally $a_1<0$ is equivalent to $a_1+a_2>\frac3{K+d+1}$: by Theorem~\ref{thm:two}(c),
$a_1+a_2-\frac3{K+d+1}=\frac{6}{(K+2)(d+2)}\big(q_{K,d}-\frac{(K-1)(d-1)}{2(K+d+1)}\big)$, because
$3+\frac{(K-1)(d-1)}{K+d+1}=\frac{(K+2)(d+2)}{K+d+1}$.
\end{proof}

\begin{corollary}[exact dimension factor]\label{cor:dim}\StmtDim\end{corollary}
\begin{corollary}[rank-one per-sample noise]\label{cor:rank1}\StmtRankOne\end{corollary}
\begin{corollary}[what is seen, what is annihilated; no domination]\label{cor:compare}\StmtCompare\end{corollary}
(The proofs of the three corollaries are at the end of the block following Theorem~\ref{thm:two}.)

\section{Location for discrete noise}\label{sec:discrete}

\begin{definition}[cube expansion]\label{def:class}\StmtClass\end{definition}
\begin{assumption}[bounded noise, weights]\label{ass:B}\StmtAssB\end{assumption}
% Module: Edgeworth comparison without Cramer's condition.
% Labels: lem:test, lem:ball, lem:cube, lem:near, prop:sums, lem:locate
\begin{lemma}[test functions]\label{lem:test}
Let $X\in\mathcal X(\varepsilon)$ with comparison density $\psie$, and $Q_j=2aj+[-a,a]^N$, $j\in\Z^N$, $a=c_\ast/\varepsilon$.
For every $f\in C_c^\infty(E)$ and every $x\in E$,
\begin{equation}\label{eq:test}
\Big|\E f(x+X)-\int f(x+y)\psie(y)\,dy\Big|\le C\Big(\varepsilon^3\|f\|_{L^1}+\sum_{j\ne0}\sup_{t\in Q_j}|\widehat f(t)|\Big),
\end{equation}
where $\widehat f(t)=\int f(z)e^{-i\inner tz}dz$. Consequently, for every integer $m>N$,
the right-hand side is at most $C_m(\varepsilon^3\|f\|_{L^1}+\varepsilon^m\|D^mf\|_{L^1})$.
\end{lemma}
\begin{proof}
By Fourier inversion, $\E f(x+X)=(2\pi)^{-N}\int\widehat f(t)e^{i\inner tx}\chi_X(t)\,dt$ and likewise for $\psie$ with
$\widehat\psie$ in place of $\chi_X$ ($\widehat f$ is a Schwartz function, $\chi_X$ and $\widehat\psie$ are bounded). On
$Q_0$, $|\widehat f|\le\|f\|_{L^1}$ and \textup{(E2)} bound the difference by $C_\ast\varepsilon^3\|f\|_{L^1}$. Outside
$Q_0$, $|t|\ge c_\ast/\varepsilon$ and $|\widehat\psie(t)|\le C(1+|t|^6)e^{-\lambda|t|^2/2}$, so
$\int_{E\setminus Q_0}|\widehat f||\widehat\psie|\le C\|f\|_{L^1}e^{-c/\varepsilon^2}$; and
$\int_{Q_j}|\widehat f||\chi_X|\le C_\ast\sup_{Q_j}|\widehat f|$ by \textup{(E1)}. For the last claim, $|t|\ge a|j|_\infty$
on $Q_j$ and $|\widehat f(t)|\le N^{m/2}|t|^{-m}\|D^mf\|_{L^1}$ (integrate by parts $m$ times in a coordinate with
$|t_k|\ge|t|/\sqrt N$); $\sum_{j\ne0}|j|_\infty^{-m}<\infty$ for $m>N$.
\end{proof}

\begin{lemma}[small balls and thin neighbourhoods of cones]\label{lem:ball}
Let $X\in\mathcal X(\varepsilon)$. \textup{(i)} $\Prob(X\in B(y,r))\le C\max(r,\varepsilon)^N$ for all $y\in E$, $r>0$.
\textup{(ii)} Let $V\subset E$ be a closed cone with $\vol\{y\in B(0,1):\dist(y,V)\le r\}\le C_Vr^c$ for $r\le1$. Then for
all $x\in E$, $R\ge1$ and $0<r\le R$,
$\Prob\big(|x+X|\le R,\ \dist(x+X,V)\le r\big)\le C\,R^{N-c}\max(r,\varepsilon)^c$.
\end{lemma}
\begin{proof}
(i) Let $k(z)=\prod_{j\le N}(\sin(z_j/2)/(z_j/2))^2$. Then $k\ge(2/\pi)^{2N}$ on $[-\pi,\pi]^N$, and $\widehat k$ is
bounded and supported in $[-1,1]^N$. For $r\ge\varepsilon$,
$\Prob(X\in B(y,r))\le(\pi/2)^{2N}\E\,k((X-y)/r)\le Cr^N\int_{|t|_\infty\le1/r}|\chi_X(t)|\,dt$, and the cube
$[-1/r,1/r]^N$ is covered by at most $(1+\varepsilon/(c_\ast r))^N\le C$ cubes of side $2c_\ast/\varepsilon$; apply
\textup{(E1)}. For $r<\varepsilon$ use monotonicity in $r$.
(ii) Let $\varrho=\max(r,\varepsilon)$ and $\{z_i\}$ a maximal $\varrho$-separated subset of
$A=\{z\in B(0,R):\dist(z,V)\le r\}$. The balls $B(z_i,\varrho/2)$ are disjoint and contained in
$\{z\in B(0,2R):\dist(z,V)\le2\varrho\}$, whose volume is at most $C_V(2R)^{N-c}(2\varrho)^c$ by homogeneity; so there are
at most $CR^{N-c}\varrho^{c-N}$ of them. $A$ is covered by the balls $B(z_i,\varrho)$; apply (i).
\end{proof}

\begin{lemma}[where the comparison field vanishes]\label{lem:locate}
Let $\psie$ be a comparison density as in Definition~\ref{def:class} and
$\Psi(x)=\int(\nu(x+y)-\nu(y))\psie(y)\,dy$. Then $\Psi\in C^\infty$,
$\nabla\Psi(x)=\int o(x+y)\psie(y)\,dy$, the derivatives of $\nabla\Psi$ of every order are bounded on $E$, and
\[
\nabla\Psi(x)=J_\Sigma(x-\hat x)+e(x),\qquad\hat x=-\tfrac16J_\Sigma^{-1}T_\Sigma[\cK],\qquad
|e(x)|\le C_A\varepsilon^3\quad\text{for }|x|\le A\varepsilon,
\]
for every $A\ge1$, with $C_A$ depending only on $A$, $\nu$, $\lambda$, $\Lambda$, $C_\ast$; and $|\hat x|\le C\varepsilon$.
\end{lemma}
\begin{proof}
Write $\widehat r(t)=\sum_\alpha c_\alpha(it)^\alpha$ (even multi-indices, $2\le|\alpha|\le6$, $|c_\alpha|\le C\varepsilon^2$).
Then $\psie=\varphi_\Sigma+\frac16\inner{\cK}{\He^\Sigma_3}\varphi_\Sigma+\sum_\alpha c_\alpha(-\partial)^\alpha\varphi_\Sigma$, and
Lemma~\ref{lem:gauss}(2) gives
$\nabla\Psi=\Phi_\Sigma+\frac16D^3\Phi_\Sigma[\cK]+\sum_\alpha c_\alpha\partial^\alpha\Phi_\Sigma$. By
Lemma~\ref{lem:gauss}(1),(3), for $|x|\le A\varepsilon$: $\Phi_\Sigma(x)=J_\Sigma x+O(|x|^3)$;
$D^3\Phi_\Sigma(x)[\cK]=T_\Sigma[\cK]+O(|\cK||x|^2)$ because $D^4\Phi_\Sigma(0)=0$; and
$\partial^\alpha\Phi_\Sigma(x)=O(|x|)$ for even $\alpha$. With $|\cK|\le C_\ast\varepsilon$ all the remainders are $O(\varepsilon^3)$.
\end{proof}

\paragraph{Weighted sums of bounded vectors.} We now verify that the normalized weighted sums of
Assumption~\ref{ass:B} belong to the class, with no condition on the characteristic function
$\varphi(s)=\E e^{i\inner s\zeta}$.

\begin{lemma}[mass of the characteristic function in cubes]\label{lem:cube}
Under Assumption~\ref{ass:B} there are $c_1,C_1>0$, depending only on $N$, $M$ and $\lambda$, such that if $\ew\le\frac12$,
\[
\int_{t_0+[-a,a]^N}\big|\E e^{i\inner t{X_w}}\big|\,dt\ \le\ C_1\qquad\text{for every }t_0\in E,\qquad a=\frac{c_1}{\ew}.
\]
\end{lemma}
\begin{proof}
With $\zeta'$ an independent copy of $\zeta$,
$|\varphi(u)|^2=\E\cos\inner u{\zeta-\zeta'}\le1-\inner u{\Sigma u}+\frac23M^4|u|^4\le e^{-\lambda|u|^2/2}$ for
$|u|\le u_0:=(3\lambda/4)^{1/2}M^{-2}$ (use $\cos\theta\le1-\theta^2/2+\theta^4/24$ and $|\zeta-\zeta'|\le2M$). Since
$\max_kw_k^2\le\rho_2/4$, the indices can be split into two sets $A,B$ with $\sum_{k\in A}w_k^2\ge\rho_2/4$ and
$\sum_{k\in B}w_k^2\ge\rho_2/4$: fill $A$ until its sum first reaches $\rho_2/4$; it is then at most $\rho_2/2$. With
$s=t/\sqrt{\rho_2}$,
\begin{gather*}
\big|\E e^{i\inner t{X_w}}\big|=\prod_k|\varphi(w_ks)|\le\tfrac12\big(g_A(s)+g_B(s)\big),\\
g_A(s)=\prod_{k\in A}|\varphi(w_ks)|^2=\E e^{i\inner s{Y_A}},\qquad Y_A=\sum_{k\in A}w_k(\zeta_k-\zeta'_k)
\end{gather*}
(the series converges in $L^2$ when there are infinitely many weights). Let $\Delta(s)=\prod_{j\le N}(1-|s_j|/\alpha)_+$
with $\alpha=u_0/(\sqrt N\max_kw_k)$; its Fourier transform
$\check\Delta(y)=\int\Delta(s)e^{i\inner sy}ds=\prod_j\alpha\big(\frac{\sin(\alpha y_j/2)}{\alpha y_j/2}\big)^2$ is
nonnegative. Hence, for every $s_0$, by Fubini,
\begin{align*}
\int g_A(s)\Delta(s-s_0)\,ds&=\E\big[e^{i\inner{s_0}{Y_A}}\check\Delta(Y_A)\big]\le\E\,\check\Delta(Y_A)=\int g_A(s)\Delta(s)\,ds\\
&\le\int e^{-\lambda\rho_2|s|^2/8}ds=\Big(\frac{8\pi}{\lambda\rho_2}\Big)^{N/2},
\end{align*}
because $|w_ks|\le u_0$ on the support of $\Delta$. As $g_A\ge0$ and $\Delta(\cdot-s_0)\ge2^{-N}$ on
$s_0+[-\alpha/2,\alpha/2]^N$, we get $\int_{s_0+[-\alpha/2,\alpha/2]^N}g_A\le2^N(8\pi/\lambda\rho_2)^{N/2}$, and the same for
$g_B$. The change of variables $t=\sqrt{\rho_2}\,s$ gives the claim with $c_1=u_0/(2\sqrt N)$.
\end{proof}

The lemma says that the mass of $|\chi_{X_w}|$ near \emph{any} frequency is at most a constant multiple of its mass near
the origin. The resonances of a lattice law are the worst case, and they are harmless against test functions whose
Fourier transform is summable over a lattice of mesh $1/\ew$.

\begin{lemma}[small frequencies]\label{lem:near}
Under Assumption~\ref{ass:B} let
$\widehat\psie_w(t)=e^{-\inner t{\Sigma t}/2}\big(1+\varepsilon_3p_1(t)+\varepsilon_4p_2(t)+\frac12\varepsilon_3^2p_1(t)^2\big)$ with
$p_1(t)=\frac16K_3[(it)^{\otimes3}]$, $p_2(t)=\frac1{24}K_4[(it)^{\otimes4}]$, $\varepsilon_3=\rho_3\rho_2^{-3/2}$,
$\varepsilon_4=\rho_4\rho_2^{-2}$. There are $c_2\le c_1$ and $C_2$ such that for $|t|_\infty\le c_2/\ew$,
\[
\big|\E e^{i\inner t{X_w}}-\widehat\psie_w(t)\big|\le C_2\,\ew^3\big(|t|^5+|t|^{9}\big)e^{-\lambda|t|^2/4}.
\]
\end{lemma}
\begin{proof}
$|\varphi(u)-1|\le M^2|u|^2/2$, so on $|u|\le1/M$ the function $\log\varphi$ is $C^\infty$ with derivatives bounded in
terms of $M$, and $\log\varphi(u)=-\frac12\inner u{\Sigma u}+\frac16K_3[(iu)^{\otimes3}]+\frac1{24}K_4[(iu)^{\otimes4}]+\varrho(u)$,
$|\varrho(u)|\le C|u|^5$. Summing over $u=w_kt/\sqrt{\rho_2}$ (all of norm at most $\ew|t|$) and using
$\rho_p\rho_2^{-p/2}\le\ew^{p-2}$,
\[
\log\E e^{i\inner t{X_w}}=-\tfrac12\inner t{\Sigma t}+A,\qquad A=\varepsilon_3p_1+\varepsilon_4p_2+R,\qquad
|R|\le C\ew^3|t|^5 .
\]
If $\ew|t|\le\sqrt Nc_2$ with $c_2$ small, then $|A|\le\lambda|t|^2/4$. Now $|e^A-1-A-\frac12A^2|\le\frac16|A|^3e^{|A|}$,
$|A|^3\le C\ew^3|t|^9$, and
$A+\frac12A^2=\varepsilon_3p_1+\varepsilon_4p_2+\frac12\varepsilon_3^2p_1^2+O(\ew^3(|t|^5+|t|^8))$. Multiply by
$e^{-\inner t{\Sigma t}/2}\le e^{-\lambda|t|^2/2}$.
\end{proof}

\begin{proposition}[weighted sums belong to the class]\label{prop:sums}
Under Assumption~\ref{ass:B} there is $\varepsilon_0>0$ such that for every weight sequence with $\ew\le\varepsilon_0$,
$X_w\in\mathcal X(\ew)$ with covariance $\Sigma$, cubic tensor $\cK=\varepsilon_3K_3$, even part
$\widehat r=\varepsilon_4p_2+\frac12\varepsilon_3^2p_1^2$, and constants $C_\ast,c_\ast$ depending only on $N$, $M$,
$\lambda$. Here $\varepsilon_3=\rho_3\rho_2^{-3/2}\le\ew$ and $\varepsilon_4=\rho_4\rho_2^{-2}\le\ew^2$.
\end{proposition}
\begin{proof}
\textup{(E1)} is Lemma~\ref{lem:cube} (cubes of half-side $c_2/\ew\le c_1/\ew$ are contained in cubes of half-side
$c_1/\ew$). \textup{(E2)}: integrate the bound of Lemma~\ref{lem:near} over $|t|_\infty\le c_2/\ew$. \textup{(E3)}:
Hoeffding's inequality for each coordinate of $X_w=\sum_k(w_k/\sqrt{\rho_2})\zeta_k$, since
$\sum_kw_k^2/\rho_2=1$, gives $\Prob(|X_w|>u)\le2Ne^{-u^2/(2NM^2)}$. The bounds on $\cK$ and $\widehat r$ follow from
$|K_3|\le M^3$, $|K_4|\le CM^4$.
\end{proof}


\subsection{Mechanism I: singular sets of codimension at least two}
\begin{assumption}[singular set of codimension $c$]\label{ass:R}\StmtAssR\end{assumption}
\begin{theorem}[codimension principle]\label{thm:codim}\StmtCodim\end{theorem}
% Module: mechanism I, singular sets of codimension >= 2.
% Labels: lem:examples, prop:field, proof of thm:codim
\begin{lemma}[examples]\label{lem:examples}
Assumption~\ref{ass:R} holds for: \textup{(a)} the Euclidean norm on $\R^N$, with $V=\{0\}$, $c=N$; \textup{(b)}
block norms $\nu(y)=\sum_g|y_g|$ over a partition of the coordinates into blocks of sizes $n_g$, with
$V=\bigcup_g\{y_g=0\}$, $c=\min_gn_g$; \textup{(c)} the nuclear norm on $\R^{K\times d}$ with $K<d$ (transpose if $K>d$),
with $o=\polar$, $V=\{\rank<K\}$, $\dist(\cdot,V)=\sigma_K$ and $c=d-K+1$. For the $\ell_1$ norm (sign descent) and for
the nuclear norm on square matrices it holds with $c=1$ only.
\end{lemma}
\begin{proof}
(a), (b): $o(y)=(y_g/|y_g|)_g$ is smooth off $V$ and $0$-homogeneous in each block, so $|D^ko(y)|\le C_k\min_g|y_g|^{-k}$
and $\dist(y,V)=\min_g|y_g|$; the volume bound is elementary. (c) For $MM^\top\succ0$,
\[
\polar(M)=(MM^\top)^{-1/2}M=\frac2\pi\int_0^\infty G_\tau(M)\,d\tau,\qquad G_\tau(M)=A^{-1}M,\quad A=MM^\top+\tau^2I .
\]
With $s=(\sigma_K(M)^2+\tau^2)^{1/2}$ we have $\|A^{-1/2}\|\le1/s$ and $\|A^{-1/2}M\|\le1$ in operator norm. Since
$D(A^{-1})[H]=-A^{-1}(HM^\top+MH^\top)A^{-1}$, the derivative $D^kG_\tau(M)[H,\dots,H]$ is a sum of at most $C_k$ words
$A^{-1}X_1A^{-1}X_2\cdots A^{-1}X_jA^{-1}Y$ with $X_i\in\{HM^\top,MH^\top,HH^\top\}$, $Y\in\{M,H\}$, and $k$ letters $H$
in total. Write each $A^{-1}=A^{-1/2}A^{-1/2}$: there are $2(j+1)$ halves; each letter $M$ or $M^\top$ can be grouped
with an adjacent half of its own into a factor of norm at most $1$, and if $\ell$ is the number of letters $M,M^\top$
then $k+\ell=2j+1$, so exactly $2(j+1)-\ell=k+1$ halves remain, each of norm at most $1/s$. Hence
$\|D^kG_\tau(M)\|\le C_k(\sigma_K^2+\tau^2)^{-(k+1)/2}$ and, integrating in $\tau$ for $k\ge1$,
$\|D^k\polar(M)\|\le C_k'\sigma_K(M)^{-k}$ (the case $k=0$ is trivial). The distance from $M$ to $V$ is $\sigma_K(M)$
(Eckart--Young). Finally, with rows $\varrho_j$ of $M$,
$\sigma_K(M)\ge K^{-1/2}\min_j\dist(\varrho_j,\operatorname{span}\{\varrho_l\}_{l\ne j})$, so $\sigma_K(M)\le r$ forces some
row to lie within $\sqrt Kr$ of the span of the others, a subspace of codimension $d-K+1$ in $\R^d$; given the other
rows this has volume at most $C r^{d-K+1}$ in the unit ball. The last sentence: $o$ jumps across a hypersurface.
\end{proof}

\begin{proposition}[mean field, codimension $\ge2$]\label{prop:field}
Let Assumption~\ref{ass:R} hold with $c\ge2$ and let $X\in\mathcal X(\varepsilon)$, $\varepsilon\le\varepsilon_0$, with
comparison density $\psie$. Then, with $\Le=\log(e/\varepsilon)$,
\[
\sup_{|x|\le1}\Big|\E\,o(x+X)-\int o(x+y)\psie(y)\,dy\Big|\le C\big(\varepsilon^{c}+\varepsilon^3\big)\Le^{N/2}.
\]
\end{proposition}
\begin{proof}
Choose $R=C'\sqrt{\Le}$ with $C'$ so large that $\Prob(|x+X|>R)\le\varepsilon^3$ for $|x|\le1$ (by \textup{(E3)}) and
$\int_{|x+y|>R}|\psie(y)|dy\le\varepsilon^3$. Let $\chi_R=\chi(\cdot/R)$ with $\chi\in C_c^\infty(B(0,2))$, $\chi=1$ on
$B(0,1)$, and let $\Theta\in C^\infty(E)$, $0\le\Theta\le1$, with $\Theta=1$ on $\{\dist(\cdot,V)\le\varepsilon\}$,
$\Theta=0$ on $\{\dist(\cdot,V)\ge2\varepsilon\}$, $|D^k\Theta|\le C_k\varepsilon^{-k}$ (mollify the indicator of
$\{\dist\le\frac32\varepsilon\}$ at scale $\varepsilon/4$). Write $o=F_1+F_2+F_3$ with $F_1=o(1-\Theta)\chi_R$,
$F_2=o\,\Theta\chi_R$, $F_3=o(1-\chi_R)$.

$F_3$: both $\E|F_3(x+X)|$ and $\int|F_3(x+y)||\psie(y)|dy$ are at most $C\varepsilon^3$.

$F_2$: $\E|F_2(x+X)|\le\sup|o|\,\Prob(|x+X|\le2R,\ \dist(x+X,V)\le2\varepsilon)\le CR^{N-c}\varepsilon^c$ by
Lemma~\ref{lem:ball}(ii), and $\int|F_2(x+y)||\psie(y)|dy\le C\vol\{z\in B(0,2R):\dist(z,V)\le2\varepsilon\}\le CR^{N-c}\varepsilon^c$.

$F_1\in C_c^\infty(E)$: on its support $\dist(z,V)\ge\varepsilon$, and by Leibniz' rule and Assumption~\ref{ass:R},
$|D^mF_1(z)|\le C_m(\dist(z,V)^{-m}+1)$. Summing over the layers $\{2^i\varepsilon\le\dist(\cdot,V)<2^{i+1}\varepsilon\}$,
$i\ge0$, inside $B(0,2R)$,
\[
\|D^mF_1\|_{L^1}\le C\sum_{i\ge0}(2^i\varepsilon)^{-m}R^{N-c}(2^{i+1}\varepsilon)^c+CR^N\le C\big(R^{N-c}\varepsilon^{c-m}+R^N\big)
\]
for $m=N+1>c$.
Lemma~\ref{lem:test} then gives an error $C(\varepsilon^3R^N+\varepsilon^cR^{N-c}+\varepsilon^mR^N)$ for $F_1$ (apply it to
each coordinate). Summing the three contributions proves the claim.
\end{proof}

\begin{proof}[Proof of Theorem~\ref{thm:codim}]
Let $H(x)=\E[\nu(x+X)-\nu(X)]$, a convex function with $\E\,o(x+X)\in\partial H(x)$. By Proposition~\ref{prop:field}
and Lemma~\ref{lem:locate}, uniformly for $|x|\le A\varepsilon$,
\[
\E\,o(x+X)=J_\Sigma(x-\hat x)+e_1(x),\qquad|e_1(x)|\le C_1\varepsilon^{\min(c,3)}\Le^{N/2}.
\]
Let $\varrho=2C_1\varepsilon^{\min(c,3)}\Le^{N/2}/\lambda_{\min}(J_\Sigma)$, which is at most $\varepsilon$ for
$\varepsilon$ small, and take $A$ with $|\hat x|+\varepsilon\le A\varepsilon$. For $|u|=1$,
$\inner{\E\,o(\hat x+\varrho u+X)}{u}\ge\lambda_{\min}(J_\Sigma)\varrho-C_1\varepsilon^{\min(c,3)}\Le^{N/2}>0$. If a minimizer
$x'$ of $H$ had $|x'-\hat x|>\varrho$, put $u=(x'-\hat x)/|x'-\hat x|$ and $y=\hat x+\varrho u$: the subgradient
inequality gives $H(x')\ge H(y)+(|x'-\hat x|-\varrho)\inner{\E\,o(y+X)}u>H(y)$, a contradiction. Hence
$\arg\min H\subset B(\hat x,\varrho)$, which is the first claim. For weighted sums, Proposition~\ref{prop:sums} gives
$X_w\in\mathcal X(\ew)$ with $\cK=\varepsilon_3K_3$, so $\hat x=\varepsilon_3\gs$; and
$\psi_{\xi_w}(G)=\sqrt{\rho_2}\,H(G/\sqrt{\rho_2})$ with $\sqrt{\rho_2}\,\varepsilon_3=\rho_3/\rho_2$.
\end{proof}


\subsection{Mechanism II: Fourier decay of the Laplacian of the norm}
\begin{assumption}[decay]\label{ass:L}\StmtAssL\end{assumption}
\begin{theorem}[curvature and symmetry]\label{thm:laplace}\StmtLaplace\end{theorem}
% Module: mechanism II, Fourier decay of the Laplacian of the norm.
% Labels: prop:value, lem:upgrade, proof of thm:laplace, lem:orbit, lem:matrices, lem:twobytwo
Throughout this part $\chi\in C_c^\infty(B(0,2))$ is a fixed radial function with $0\le\chi\le1$ and $\chi=1$ on $B(0,1)$,
$\mu=\Delta\nu$ and $\lambda_\nu=\Delta(\chi\nu)$ in the sense of distributions.

\begin{lemma}[the Laplacian of a norm]\label{lem:lapmeasure}
$\mu$ is a nonnegative Radon measure with $\mu(B(0,s))=s^{N-1}\mu(B(0,1))$, and
$\lambda_\nu=\chi\mu+(2\inner{\nabla\chi}{\nabla\nu}+\nu\Delta\chi)\,dz$ is a finite signed measure supported in
$\overline{B(0,2)}$ with $|\lambda_\nu|(B(0,s))\le Cs^{N-1}$ for $s>0$. With $F_R(z)=\chi(z/R)\nu(z)$, $R\ge1$,
\[
\widehat{F_R}(t)=R^{N+1}\widehat{F_1}(Rt),\qquad|t|^2\widehat{F_1}(t)=-\widehat{\lambda_\nu}(t),\qquad
\|F_R\|_{L^1}\le CR^{N+1}.
\]
Consequently, under Assumption~\ref{ass:L}, $|\widehat{F_R}(t)|\le CR^{N+1}(1+R|t|)^{-2-\kappa}$.
\end{lemma}
\begin{proof}
A convex function has a distributional Hessian which is a nonnegative matrix-valued measure (mollify: the Hessians of
$\nu*h_s$ are nonnegative and converge weakly); its trace is $\mu$. Since $\nu(sz)=s\nu(z)$, the distribution $\mu$ is
homogeneous of degree $-1$, which gives the scaling. The product rule gives the formula for $\lambda_\nu$; the two
last terms are bounded and supported in $\{1\le|z|\le2\}$. The identities for $F_R$ are the dilation rule
$F_R(z)=RF_1(z/R)$ and $\widehat{\Delta f}=-|t|^2\widehat f$. Finally $|\widehat{F_1}(t)|\le\|F_1\|_{L^1}$ for $|t|\le1$
and $|\widehat{F_1}(t)|=|t|^{-2}|\widehat{\lambda_\nu}(t)|$ for $|t|\ge1$.
\end{proof}

\begin{proposition}[comparison of potentials]\label{prop:value}
Let Assumption~\ref{ass:L} hold and $X\in\mathcal X(\varepsilon)$, $\varepsilon\le\varepsilon_0$, with comparison density
$\psie$. Let $H(x)=\E[\nu(x+X)-\nu(X)]$ and $\Psi(x)=\int(\nu(x+y)-\nu(y))\psie(y)\,dy$. Then, with
$q=N/(N-\kappa)$ and $\Le=\log(e/\varepsilon)$,
\[
\sup_{|x|\le2}|H(x)-\Psi(x)|\le C\,\varepsilon^{2q}\,\Le^{(N+1)/2}.
\]
\end{proposition}
\begin{proof}
Let $h\in C_c^\infty(B(0,1))$ be even, nonnegative, of integral one, $h_s=s^{-N}h(\cdot/s)$, $\delta=\varepsilon^q$
(so $\varepsilon\ge\delta$) and $R=C'\sqrt{\Le}$.

\emph{Step 1: smoothed, truncated norm.} Let $f=F_R*h_s$ with $s\in\{\delta,2\delta\}$. Then $f\in C_c^\infty$,
$\|f\|_{L^1}\le CR^{N+1}$ and $|\widehat f(t)|=|\widehat{F_R}(t)||\widehat h(st)|\le C_DR^{N+1}(1+R|t|)^{-2-\kappa}(1+s|t|)^{-D}$
for any $D$, by Lemma~\ref{lem:lapmeasure}. On $Q_j$, $j\ne0$, we have $|t|\ge c_\ast|j|_\infty/\varepsilon$, so by
Lemma~\ref{lem:test}
\[
\Big|\E f(x+X)-\!\int\! f(x+y)\psie(y)dy\Big|\le C\varepsilon^3R^{N+1}+CR^{N-1-\kappa}\varepsilon^{2+\kappa}\sum_{j\ne0}\frac{(1+s|j|/\varepsilon)^{-D}}{|j|^{2+\kappa}}.
\]
Since $2+\kappa<N$, the sum over $|j|\le\varepsilon/s$ is at most $C(\varepsilon/s)^{N-2-\kappa}$, and so is the sum over
$|j|>\varepsilon/s$ when $D>N$. The second term is therefore at most
$CR^{N-1-\kappa}\varepsilon^{N}s^{-(N-2-\kappa)}\le CR^{N-1-\kappa}\delta^{N-\kappa}\delta^{-(N-2-\kappa)}=CR^{N-1-\kappa}\delta^2$,
because $\varepsilon^N=\delta^{N-\kappa}$. Also $\varepsilon^3\le\delta^2$, since $2q\le3$ is equivalent to $N\ge3\kappa$.

\emph{Step 2: removing the truncation.} For $|x|\le2$,
$|(\nu*h_s-F_R*h_s)(z)|\le C(1+|z|)\mathbf 1\{|z|\ge R-1\}$; by \textup{(E3)} and the Gaussian tails of $\psie$, both
$\E$ and $\int\cdot\,|\psie|$ of this function at $x+\cdot$ are at most $\varepsilon^3$ if $C'$ is large. Hence
\begin{align*}
\E(\nu*h_s)(x+X)&=\int(\nu*h_s)(x+y)\psie(y)\,dy+O(R^{N+1}\delta^2)\\&=(\Pi*h_s)(x)+O(R^{N+1}\delta^2),
\end{align*}
where $\Pi(x)=\int\nu(x+y)\psie(y)dy$. $\Pi$ is smooth with bounded second derivatives (Lemma~\ref{lem:locate}) and
$h_s$ is even, so $(\Pi*h_s)(x)=\Pi(x)+O(s^2)$.

\emph{Step 3: convexity removes the smoothing.} For every $z\in E$,
\begin{equation}\label{eq:sandwich}
2(\nu*h_\delta)(z)-(\nu*h_{2\delta})(z)\ \le\ \nu(z)\ \le\ (\nu*h_\delta)(z).
\end{equation}
The right inequality is Jensen's ($h_\delta$ is even with integral one). The left one: $2\nu(z-u)\le\nu(z)+\nu(z-2u)$
by convexity; integrate against $h_\delta(u)\,du$ and note $\int\nu(z-2u)h_\delta(u)du=(\nu*h_{2\delta})(z)$. Taking the
expectation of \eqref{eq:sandwich} at $z=x+X$ under the \emph{true} law of $X$, and using Step 2 for both
$s=\delta$ and $s=2\delta$,
$\E\nu(x+X)=\Pi(x)+O(R^{N+1}\delta^2)$ for $|x|\le2$. Subtract the same identity at $x=0$.
\end{proof}

Step 3 is the point where atoms cease to matter: no smoothness of the law of $X$ and no positivity of the signed
density $\psie$ are used, only the convexity of $\nu$.

\begin{lemma}[from values to subgradients]\label{lem:upgrade}
Let $H$ be convex on $B(0,2)\subset E$, $\Psi\in C^2$ with $\|D^2\Psi\|\le C_2$ there, and
$\sup_{B(0,2)}|H-\Psi|\le\eta\le1$. Then for every $|x|\le1$ and every $v\in\partial H(x)$,
$|v-\nabla\Psi(x)|\le(2+C_2/2)\sqrt\eta$.
\end{lemma}
\begin{proof}
For $|u|=1$ and $d=\sqrt\eta$: $\inner vu\le(H(x+du)-H(x))/d\le(\Psi(x+du)-\Psi(x))/d+2\eta/d\le
\inner{\nabla\Psi(x)}u+C_2d/2+2\eta/d$. Apply this to $u$ and $-u$.
\end{proof}

\begin{proof}[Proof of Theorem~\ref{thm:laplace}]
By Proposition~\ref{prop:value} and Lemma~\ref{lem:upgrade}, with $d_\varepsilon=\varepsilon^q\Le^{(N+1)/4}$,
\begin{equation}\label{eq:allsub}
\sup_{|x|\le1}\ \sup_{v\in\partial H(x)}|v-\nabla\Psi(x)|\le C_1d_\varepsilon .
\end{equation}
By Lemma~\ref{lem:locate}, $\nabla\Psi(x)=J_\Sigma(x-\hat x)+O(\varepsilon^3)$ for $|x|\le A\varepsilon$, and
$\varepsilon^3\le d_\varepsilon\le\varepsilon$ for $\varepsilon$ small because $1<q\le\frac32$. Let
$\varrho=3C_1d_\varepsilon/\lambda_{\min}(J_\Sigma)$. For $|u|=1$, every $v\in\partial H(\hat x+\varrho u)$ has
$\inner vu\ge\lambda_{\min}(J_\Sigma)\varrho-2C_1d_\varepsilon>0$. A minimizer $x'$ of $H$ with $|x'-\hat x|>\varrho$ would
contradict the subgradient inequality at $y=\hat x+\varrho u$, $u=(x'-\hat x)/|x'-\hat x|$:
$H(x')\ge H(y)+(|x'-\hat x|-\varrho)\inner vu>H(y)$. Minimizers exist since $H(x)\ge\nu(x)-2\E\nu(X)$. The statement for
weighted sums follows from Proposition~\ref{prop:sums} as in the proof of Theorem~\ref{thm:codim}.
\end{proof}

\begin{lemma}[symmetry gives decay]\label{lem:orbit}
Let $N\ge2$ and let $\mathcal G\subset O(E)$ be a compact group, with normalized Haar measure $dg$, such that $\nu\circ g=\nu$
for $g\in\mathcal G$ and, for some $0<\kappa\le1$,
\begin{equation}\label{eq:orbital}
\sup_{|e|=|y|=1}\Big|\int_{\mathcal G}e^{iu\inner e{gy}}dg\Big|\le C(1+u)^{-\kappa}\qquad(u\ge0).
\end{equation}
Then $|\widehat{\lambda_\nu}(t)|\le C'(1+|t|)^{-\kappa}$.
\end{lemma}
\begin{proof}
$\nu$, $\chi$ and the Laplacian are $\mathcal G$-invariant, hence so is $\lambda_\nu$, and
$\widehat{\lambda_\nu}(t)=\int\!\int_{\mathcal G}e^{-i\inner t{gz}}dg\,d\lambda_\nu(z)$. By \eqref{eq:orbital} and
Lemma~\ref{lem:lapmeasure},
\[
|\widehat{\lambda_\nu}(t)|\le C\int(1+|t||z|)^{-\kappa}d|\lambda_\nu|(z)\le C\sum_{i\ge-1}(1+|t|2^{-i-1})^{-\kappa}2^{-i(N-1)}
\le C'(1+|t|)^{-\kappa},
\]
summing over the annuli $\{2^{-i-1}<|z|\le2^{-i}\}$ and using $N-1\ge\kappa$.
\end{proof}

\begin{lemma}[left and right rotations]\label{lem:matrices}
Let $k,l\ge3$. The group $\mathcal G=\SO(k)\times\SO(l)$ acting on $\R^{k\times l}$ by $y\mapsto Uy V^\top$ satisfies
\eqref{eq:orbital} with $\kappa=\frac12$. Consequently every norm on $\R^{k\times l}$ that is invariant under left and
right rotations (the nuclear norm, the operator norm, all Schatten and Ky Fan norms) satisfies
Assumption~\ref{ass:L} with $\kappa=\frac12$.
\end{lemma}
\begin{proof}
Fix unit $e,y$ and a point $g_0=(U_0,V_0)$, and consider the phase $\phi(g)=\inner e{gy}$ in the chart
$(A,B)\mapsto(U_0e^{A},V_0e^{B})$, $A\in\mathfrak{so}(k)$, $B\in\mathfrak{so}(l)$. Replacing $e$ by $U_0^\top eV_0$ we may
take $g_0$ to be the identity. The first and second derivatives of $\phi$ at the origin of the chart are the numbers
$\inner e{\mathcal Dy}$, where $\mathcal D$ ranges over the linear span $\mathfrak D$ of the maps
\[
y\mapsto Ay,\quad y\mapsto yB,\quad y\mapsto A_1A_2y,\quad y\mapsto yB_1B_2,\quad y\mapsto AyB,
\]
with $A,A_i\in\mathfrak{so}(k)$ and $B,B_i\in\mathfrak{so}(l)$.
We claim that $\{\mathcal Dy:\mathcal D\in\mathfrak D\}=\R^{k\times l}$ for every $y\ne0$. With the matrix units
$E_{ij}$ and $A_{ij}=E_{ij}-E_{ji}$, one has $A_{ih}A_{hj}=-E_{ij}$ for distinct $i,h,j$ and
$A_{ij}^2=-(E_{ii}+E_{jj})$; as the size is at least $3$, products of two skew matrices span all square matrices. So the
span contains $Cy$ and $yD$ for all $C\in\R^{k\times k}$, $D\in\R^{l\times l}$. Choose orthonormal bases in which
$y=\sum_{i\le r}s_iE_{ii}$, $s_i>0$ (conjugation preserves the skew matrices). Then $Cy$ and $yD$ give every $E_{pq}$
with $q\le r$ or $p\le r$; and for $p,q>r$ and any $i\le r$, $A_{pi}\,y\,A_{iq}=s_iE_{pi}A_{iq}=s_iE_{pq}$. This proves the
claim, so for every unit $e$ some derivative of $\phi$ of order one or two at $g_0$ is nonzero. By continuity and
compactness of $\{|e|=1\}\times\{|y|=1\}\times\mathcal G$ there are $\eta>0$ and a finite atlas such that at every point
some coordinate derivative of order one or two has modulus at least $\eta$ on a coordinate ball of fixed radius, while
the $C^3$ norms of $\phi$ in the charts are bounded uniformly in $(e,y)$. With a partition of unity subordinate to the
atlas, the van der Corput estimate in several variables \citep[Ch.~VIII, \S2.2, Prop.~5]{stein}, whose constant depends
only on the lower bound $\eta$, the order of the derivative and the $C^{3}$ norm of the phase, gives
$|\int_{\mathcal G}e^{iu\phi}dg|\le Cu^{-1/2}$ uniformly. The consequence is Lemma~\ref{lem:orbit}.
\end{proof}

The exponent $\frac12$ of Lemma~\ref{lem:matrices} cannot be improved in general: for $k=l=3$ and $e=y=I/\sqrt3$ the
orbital integral is $\E e^{iu\tr(U)/3}$ for a Haar rotation $U$, and $\tr U=1+2\cos\theta$ has a density with a
non-degenerate critical value, so the integral is of exact order $u^{-1/2}$. For size $2$ the lemma fails
($\SO(2)\times\SO(2)$ fixes the two planes of $2\times2$ matrices described next, and for $e$ in one and $y$ in the other the
orbital integral is identically $1$), but Assumption~\ref{ass:L} still holds:

\begin{lemma}[$2\times2$ matrices]\label{lem:twobytwo}
For $M=\begin{psmallmatrix}p&q\\r&s\end{psmallmatrix}$ let $y_+=\frac1{\sqrt2}(p+s,\,q-r)$ and
$y_-=\frac1{\sqrt2}(p-s,\,q+r)$. Then $M\mapsto(y_+,y_-)$ is an isometry of $(\R^{2\times2},\nF\cdot)$ onto
$\R^2\times\R^2$, $\det M=\frac12(|y_+|^2-|y_-|^2)$,
\[
\nnuc M=\sqrt2\max(|y_+|,|y_-|),\qquad
\polar(M)=\begin{cases}\sqrt2\,\big(y_+/|y_+|,\,0\big)&\text{if }|y_+|>|y_-|,\\[2pt]
\sqrt2\,\big(0,\,y_-/|y_-|\big)&\text{if }|y_-|>|y_+|,\end{cases}
\]
and the nuclear norm on $\R^{2\times2}$ satisfies Assumption~\ref{ass:L} with $N=4$, $\kappa=\frac12$.
\end{lemma}
\begin{proof}
The first three assertions are direct computations: $|y_+|^2+|y_-|^2=p^2+q^2+r^2+s^2$,
$|y_+|^2-|y_-|^2=2(ps-qr)$, and $\nnuc M^2=\nF M^2+2|\det M|=2\max(|y_+|^2,|y_-|^2)$; the formula for the polar factor
is the gradient of $\nnuc\cdot$ on the open sets $\Omega_\pm=\{\pm\det>0\}$. Let $V=\{|y_+|=|y_-|\}$; $V\setminus\{0\}$ is a
smooth hypersurface with unit normal $n=(y_+,-y_-)/|y|$ pointing into $\Omega_+$. On $\Omega_\pm$,
$\nu=\sqrt2|y_\pm|$ is smooth with $\Delta\nu=\sqrt2/|y_\pm|\le2/|y|$, and across $V$ the gradient jumps by
$\sqrt2(\hat y_+,-\hat y_-)$, so that (Green's formula on $E\setminus B(0,\varrho)$, $\varrho\to0$)
\[
\mu=\Delta\nu=\sqrt2\Big(\frac{\mathbf 1_{\Omega_+}}{|y_+|}+\frac{\mathbf 1_{\Omega_-}}{|y_-|}\Big)dy+2\,dS_V .
\]
The rotations of $y_+$ alone and of $y_-$ alone are isometries of $E$ preserving $\nu$, $\chi$, $V$ and $dS_V$, hence
$\lambda_\nu$. Averaging over both,
$\widehat{\lambda_\nu}(t)=\int J_0(|t_+||y_+|)\,J_0(|t_-||y_-|)\,d\lambda_\nu(y)$ with $J_0$ the Bessel function,
$|J_0(u)|\le C(1+u)^{-1/2}$. Choose $\varsigma\in\{+,-\}$ with $|t_\varsigma|\ge|t|/\sqrt2$ and let $\varsigma'$ be the other
sign; then $|\widehat{\lambda_\nu}(t)|\le C\int(1+|t||y_\varsigma|)^{-1/2}d|\lambda_\nu|(y)$, where
$|\lambda_\nu|\le C\mathbf 1_{\{|y|\le2\}}\big[(1+|y_+|^{-1}+|y_-|^{-1})dy+dS_V\big]$. The four contributions are at most
$C|t|^{-1/2}$: the bounded density and the term $|y_{\varsigma'}|^{-1}$ because
$\int_{|y_\varsigma|\le2}(|t||y_\varsigma|)^{-1/2}d^2y_\varsigma\le C|t|^{-1/2}$ and
$\int_{|y_{\varsigma'}|\le2}(1+|y_{\varsigma'}|^{-1})d^2y_{\varsigma'}<\infty$; the term $|y_\varsigma|^{-1}$ because
$\int_{|y_\varsigma|\le2}|y_\varsigma|^{-1}(|t||y_\varsigma|)^{-1/2}d^2y_\varsigma\le C|t|^{-1/2}$; and the surface term because
$|y_\varsigma|=|y|/\sqrt2$ on $V$ and $dS_V=c\,\varrho^2d\varrho\,d\theta_+d\theta_-$ in the coordinates
$y_\pm=\frac{\varrho}{\sqrt2}(\cos\theta_\pm,\sin\theta_\pm)$. Since also $|\widehat{\lambda_\nu}|\le|\lambda_\nu|(E)<\infty$,
$|\widehat{\lambda_\nu}(t)|\le C(1+|t|)^{-1/2}$, and $2+\frac12<4$.
\end{proof}

So a Muon-type step on a $2\times2$ matrix normalizes the larger of the conformal and anticonformal parts of its
argument, and the singular set $\{\det=0\}$ is the cone over a Clifford torus.


\subsection{The nuclear norm in every size; consequences}
\begin{theorem}[orthogonalized steps and discrete noise]\label{thm:nuclear}\StmtNuclear\end{theorem}
\begin{corollary}[batch and momentum together, uniformly in $\beta$]\label{cor:mom}\StmtMomCor\end{corollary}
\begin{corollary}[finite sums]\label{cor:fs}\StmtFS\end{corollary}
% Module: proofs of thm:nuclear, cor:mom (uniform in beta), cor:fs
\begin{proof}[Proof of Theorem~\ref{thm:nuclear}]
If $K\ne d$, Lemma~\ref{lem:examples}(c) gives Assumption~\ref{ass:R} with $c=|K-d|+1\ge2$ (for $K=1$ or $d=1$ the nuclear
norm is the Euclidean norm of $\R^N$, $N\ge2$, and $c=N=|K-d|+1$), and Theorem~\ref{thm:codim} applies:
$\min(c,3)-1=2\gamma$. If $K=d=2$, Lemma~\ref{lem:twobytwo} gives Assumption~\ref{ass:L} with $N=4$, $\kappa=\frac12$; if
$K=d\ge3$, Lemma~\ref{lem:matrices} gives it with $N=K^2$, $\kappa=\frac12$; in both cases Theorem~\ref{thm:laplace}
applies and $q-1=\frac{\kappa}{N-\kappa}=2\gamma$. Finally $\sqrt{\rho_2}\,\ew=\max_kw_k$, which is $1/b$ for a batch mean.
\end{proof}

\begin{proof}[Proof of Corollary~\ref{cor:mom}]
$N_\beta=\xi_w$ for the weights $(1-\beta)\beta^k/b$, each repeated $b$ times: $\rho_p=(1-\beta)^pb^{1-p}/(1-\beta^p)$,
so $\rho_3/\rho_2=1/\beff$, $\ew^2=(1-\beta^2)/b=\theta$ and $\max_kw_k=(1-\beta)/b\le3/\beff$. Apply
\eqref{eq:codimw} or \eqref{eq:laplacew}.
\end{proof}

\begin{proof}[Proof of Corollary~\ref{cor:fs}]
Theorem~\ref{thm:nuclear} applied to $\zeta=\nabla f_I(W)-\nabla F(W)$, which is bounded, centred, with covariance
positive definite exactly when the $n$ gradients affinely span $E$. If $n=N+1$, the differences
$\nabla f_i(W)-\nabla f_1(W)$, $i\ge2$, form a basis of $E$ and generate a lattice that contains the support of
$\zeta-\zeta'$: the law of $\zeta$ is a lattice law.
\end{proof}


\subsection{The flat case: sign descent}
\begin{theorem}[sign descent with discrete noise]\label{thm:sign}\StmtSign\end{theorem}
% Module: the flat case (sign descent). Label: proof of thm:sign
\begin{proof}[Proof of Theorem~\ref{thm:sign}]
$b\,G^\st_b=-\{\text{medians of }S_b\}$. \textup{(a)} Esseen's theorem for non-lattice laws with three moments
\citep{esseen,gk}: uniformly in $x$,
$\Prob(S_b\le x\sigma\sqrt b)=\Phi(x)+\frac{\kappa_3}{6\sigma^3\sqrt b}(1-x^2)\varphi(x)+o(b^{-1/2})$, with $\Phi,\varphi$
the standard normal distribution function and density. At $x=m/(\sigma\sqrt b)$ with $m$ bounded this reads
$\frac12+\frac{\varphi(0)}{\sigma\sqrt b}(m+\gs)+o(b^{-1/2})$, and the same holds for $\Prob(S_b<m)$; so every
median $m_b$ of $S_b$ satisfies $m_b+\gs\to0$ (medians are bounded, by the same expansion at $|m|$ large).

\textup{(b)} Esseen's expansion for lattice laws \citep[\S43]{gk} reads, uniformly in $x$,
\[
\Prob(S_b\le x\sigma\sqrt b)=\Phi(x)+\frac{\kappa_3(1-x^2)}{6\sigma^3\sqrt b}\varphi(x)+\frac h{\sigma\sqrt b}\,
S_1\Big(\frac{x\sigma\sqrt b-ba}{h}\Big)\varphi(x)+o(b^{-1/2}),
\]
where $S_1(y)=\lfloor y\rfloor-y+\tfrac12$.
At a point $m\in ba+h\Z$ of the lattice carrying $S_b$, $S_1=\frac12$, and immediately to the left of the next lattice
point $S_1\to-\frac12$. Hence, for bounded $m\in ba+h\Z$,
\begin{align*}
\Prob(S_b\le m)&=\tfrac12+\tfrac{\varphi(0)}{\sigma\sqrt b}\big(m+\gs+\tfrac h2\big)+o(b^{-1/2}),\\
\Prob(S_b<m)=\Prob(S_b\le m-h)&=\tfrac12+\tfrac{\varphi(0)}{\sigma\sqrt b}\big(m+\gs-\tfrac h2\big)+o(b^{-1/2}).
\end{align*}
A lattice point $m$ is a median if and only if $\Prob(S_b<m)\le\frac12\le\Prob(S_b\le m)$, and a point off the lattice
is a median only if it lies between two consecutive lattice medians. So for $b$ large every median lies in
$[-\gs-\frac h2-\delta,-\gs+\frac h2+\delta]$, and a lattice point with $|m+\gs|\le\frac h2-\delta$ has
$\Prob(S_b<m)<\frac12<\Prob(S_b\le m)$, hence is the unique median. \textup{(i)} follows from Weyl's
equidistribution of $ba/h$ modulo $1$. \textup{(ii)}: $\Lambda_b$ depends on $b$ modulo $q$ only; if $q\ge2$ the cosets
are distinct, at most one of them has $\gs$ at the midpoint between two of its points, and the nearest points of two
different cosets differ, so there is no limit; if $q=1$, $\Lambda_b=h\Z$. \textup{(iii)}: $a=0$ and $|\gs|<h/2$.
\end{proof}


\subsection{Sampling without replacement}
\begin{theorem}[finite-population bias]\label{thm:wor}\StmtWor\end{theorem}
% Module: sampling without replacement. Labels: prop:worclass, proof of thm:wor
\begin{proposition}[the conditional law belongs to the class]\label{prop:worclass}
In the setting of Theorem~\ref{thm:wor} let $h=\sqrt v$ and $T=b\,\xi_{n,b}/h$. There is $v_0$ such that for $v\ge v_0$,
$T\in\mathcal X(1/h)$ with covariance $\Sigma_n$, cubic tensor $\cK=\frac{1-2p}{h}K_{3,n}$, and constants
$\lambda,\Lambda=M^2,C_\ast,c_\ast$ depending only on $N,M,\lambda$.
\end{proposition}
\begin{proof}
Let $I_1,\dots,I_n$ be independent Bernoulli variables of parameter $p=b/n$, $r=p(1-p)$, so $h^2=nr$, and
\[
T'=\frac1h\sum_i(I_i-p)z_i,\qquad C'=\sum_i(I_i-p)=\sum_iI_i-b\ \in\Z .
\]
Conditionally on $C'=0$ the set $\{i:I_i=1\}$ is uniform among subsets of size $b$, and since $\sum_iz_i=0$, $T'$ has
the law of $T$. With $\varphi_p(s)=\E e^{is(I-p)}$, $a_i(t)=\inner t{z_i}$, $s_i=(a_i(t)+u)/h$ and
$H_0(t,u)=\E e^{i\inner t{T'}+iuC'/h}=\prod_i\varphi_p(s_i)$, the identity
$\mathbf 1\{C'=0\}=\frac1{2\pi}\int_{-\pi}^{\pi}e^{i\theta C'}d\theta$ gives (the classical device for finite-population limit theorems, going back to \citet{erdosrenyi}; see
\citet{bs} for Edgeworth expansions in this setting)
\begin{equation}\label{eq:condcf}
\chi_T(t)=\frac{\int_{-\pi h}^{\pi h}H_0(t,u)\,du}{\int_{-\pi h}^{\pi h}H_0(0,u)\,du},\qquad
\int_{-\pi h}^{\pi h}\!H_0(0,u)\,du=2\pi h\,\Prob(C'=0)=\sqrt{2\pi}\,\big(1+O(\tfrac1v)\big),
\end{equation}
the last estimate being Stirling's formula for the central binomial probability, uniform in $p$.

\emph{\textup{(E1)}.} $|\varphi_p(s)|^2=1-2r(1-\cos s)$, hence $|H_0(t,u)|\le D(t,u):=\exp(-r\sum_i(1-\cos s_i))$. The function
$D$ is nonnegative and is the characteristic function of the random vector
$Y=\frac1h\sum_i(P_i-P_i')\,(z_i,1)\in E\times\R$, where $P_i,P_i'$ are independent Poisson variables of parameter $r/2$
(indeed $\E e^{is(P-P')}=e^{r(\cos s-1)}$). Let $\alpha=h/(M\sqrt N+1)$ and $\Delta$ the product of the $N+1$ tents
$(1-|\tau_j|/\alpha)_+$. On the support of $\Delta$, $|s_i|\le1$, so $1-\cos s_i\ge s_i^2/4$ and, because $\sum_iz_i=0$ and
$h^2=nr$,
\[
r\sum_is_i^2=\inner t{\Sigma_nt}+u^2,\qquad D(t,u)\le e^{-(\lambda|t|^2+u^2)/4}.
\]
As in the proof of Lemma~\ref{lem:cube}, positivity of $\check\Delta$ gives, for every $\tau_0\in E\times\R$,
$\int D\,\Delta(\cdot-\tau_0)=\E[e^{i\inner{\tau_0}Y}\check\Delta(Y)]\le\int D\,\Delta\le C$, hence $\int D\le2^{N+1}C$ on every
cube of half-side $\alpha/2$. A set $(t_0+[-a,a]^N)\times[-\pi h,\pi h]$ with $a=c_\ast h$ is covered by a bounded number of
such cubes, and with \eqref{eq:condcf} this gives $\int_{t_0+[-a,a]^N}|\chi_T|\le C$.

\emph{\textup{(E2)}.} Let $d=1-2p$ and $d_4=1-6r$. The function $(p,z)\mapsto\log(1+p(e^z-1))-pz$ is analytic in
$|z|\le z_0$, smooth in $p\in[0,1]$, and vanishes identically for $p=0$ and $p=1$; dividing by $p(1-p)$ leaves a family
of analytic functions bounded uniformly in $p$. Hence, uniformly in $p\in(0,1)$, for real $|s|\le z_0$,
$\log\varphi_p(s)=r\big[\frac{(is)^2}2+d\frac{(is)^3}6+d_4\frac{(is)^4}{24}+O(|s|^5)\big]$. Fix $c_2\le z_0/2$ and
$c_1$ with $M\sqrt Nc_1\le c_2/2$. For $|t|_\infty\le c_1h$ and $|u|\le c_2h$, summing over $i$ with
$A_j=\frac1n\sum_l(i(a_l(t)+u))^j$,
\[
H_0(t,u)=e^{-(\inner t{\Sigma_nt}+u^2)/2}\exp\Big(\frac{d}{6h}A_3+\frac{d_4}{24h^2}A_4+R\Big),\quad|R|\le\frac C{h^3}(|t|+|u|)^5 ,
\]
and expanding the exponential as in Lemma~\ref{lem:near} (after reducing $c_1,c_2$),
\begin{align*}
H_0(t,u)={}&e^{-(\inner t{\Sigma_nt}+u^2)/2}\Big(1+\frac{d}{6h}A_3+\frac1{h^2}\Big\{\frac{d_4}{24}A_4+\frac{d^2}{72}A_3^2\Big\}\Big)\\
&+O\Big(h^{-3}(1+|t|+|u|)^{15}e^{-c(|t|^2+u^2)}\Big).
\end{align*}
For $c_2h\le|u|\le\pi h$ and $|t|_\infty\le c_1h$, every $s_i$ is at distance at least $c_2/2$ from $2\pi\Z$, so
$|H_0|\le D\le e^{-c'h^2}$. We may therefore integrate the Gaussian polynomial over $u\in\R$ at the cost of an error
$O(h^{-3})$ after integration in $t$ over $[-c_1h,c_1h]^N$. Odd powers of $u$ integrate to zero against $e^{-u^2/2}$
and $\frac1n\sum_la_l(t)=0$, so
$\int A_3e^{-u^2/2}du=\sqrt{2\pi}\,\frac1n\sum_l(ia_l(t))^3=\sqrt{2\pi}\,K_{3,n}[(it)^{\otimes3}]$, while the term in braces
integrates to $\sqrt{2\pi}\,q_2(t)$ with $q_2$ an even polynomial of degree at most $6$ with coefficients bounded in terms
of $M$ (it involves $\frac1n\sum_la_l^2$, $\frac1n\sum_la_l^4$ and $(\frac1n\sum_la_l^3)^2$). The same computation at
$t=0$ gives the denominator of \eqref{eq:condcf} as $\sqrt{2\pi}(1+h^{-2}q_2(0)+O(h^{-3}))$. Dividing,
\[
\int_{[-c_1h,c_1h]^N}\Big|\chi_T(t)-e^{-\inner t{\Sigma_nt}/2}\Big(1+\frac{d}{6h}K_{3,n}[(it)^{\otimes3}]+\frac{q_2(t)-q_2(0)}{h^2}\Big)\Big|\,dt\le\frac C{h^3},
\]
which is \textup{(E2)} with $\widehat r=h^{-2}(q_2-q_2(0))$ and $c_\ast=c_1$.

\emph{\textup{(E3)}.} Replacing the sample by its complement changes $b\,\xi_{n,b}$ into its opposite, so we may assume
that $m_0=\min(b,n-b)\le2h^2$ points are drawn. By \citet[Theorem~4]{hoeffding}, for a coordinate functional $y$ and
real $s$, $\E e^{s\,y(b\xi_{n,b})}$ is at most the same expectation for sampling with replacement, which is at most
$e^{m_0s^2M^2/2}$. Since $T=b\,\xi_{n,b}/h$ and $m_0/h^2\le2$, the coordinates of $T$ are sub-Gaussian with variance proxy $2M^2$.
\end{proof}

\begin{proof}[Proof of Theorem~\ref{thm:wor}]
By Proposition~\ref{prop:worclass}, Theorem~\ref{thm:codim} or Theorem~\ref{thm:laplace} applies to $X=T$ with
$\varepsilon=1/h$: every minimizer $x$ of $x\mapsto\E\nu(x+T)$ satisfies
$|x-\frac{1-2p}h\gs^{(n)}|\le C\,h^{-1-2\gamma}\log^{N}(eh)$ (the constants are uniform because $\Sigma_n$ ranges in a
compact set of positive definite matrices). Since $\psi_{\xi_{n,b}}(G)=\frac hb\,\E[\nu(\frac bhG+T)-\nu(T)]$, the set
$\arg\min\psi_{\xi_{n,b}}=-\Med_\nu(\xi_{n,b})$ is $\frac hb$ times the set of these minimizers. Multiplying by $h$ gives
\eqref{eq:wor} with $1-2p$ in place of $\frac{n-2b}{n-2}$; the two differ by $\frac{2(n-2b)}{n(n-2)}=O(1/n)=O(1/v)$, and
$\gs^{(n)}$ is bounded. (The coefficient $\frac{n-2b}{n-2}$ is the one that matches the exact conditional cumulants,
$\Cov(T)=\frac n{n-1}\Sigma_n$ and $\E T^{\otimes3}=\frac{1-2p}{h}\frac{n^2}{(n-1)(n-2)}K_{3,n}$, and the exact duality
$\Med_\nu(\xi_{n,n-b})=-\frac b{n-b}\Med_\nu(\xi_{n,b})$.)
\end{proof}


\subsection{Sharpness}
\begin{proposition}[formal; confirmed numerically]\label{prop:lattice2}\StmtLattice\end{proposition}
% Module: second-order lattice term (formal derivation). Label: none
\noindent\emph{Derivation of Proposition~\ref{prop:lattice2} (formal).} In the scale of sums,
$\E\,o(m+S_b)=\sum_{\ell\in\Lambda_b}o(m+\ell)\Prob(S_b=\ell)$, and the local limit theorem gives
$\Prob(S_b=\ell)=|\cL|\varphi_{b\Sigma}(\ell)(1+b^{-1/2}P_1(\ell/\sqrt b)+\dots)$. The lattice sum of the kernel
$o(m+\cdot)$, homogeneous of degree $0$ and singular at $-m$, differs from its integral by
$\varphi_{b\Sigma}(0)\,\mathfrak E_{\Lambda_b}(m)$ (take $\tau^2=b$ in the definition) up to smaller terms; the integral is
$b^{-1/2}J_\Sigma(m-\gs)+O(b^{-3/2})$ by Lemma~\ref{lem:locate}; and $\varphi_{b\Sigma}(0)=(2\pi b\sqrt{\det\Sigma})^{-1}$.
Solving for the zero $m$ gives the formula. A proof requires the lattice local expansion with a remainder summable
against a non-integrable kernel, which we have not written; the statement is supported by the exact computations of
Table~\ref{tab:exact}.


\section{Verification record}\label{sec:verif}

Every statement of this document was checked in three independent ways where applicable: (a) a line-by-line review by
a reader who was given the proof and asked to break it (reports in \texttt{verification/} and in the note folders);
(b) exact or high-precision numerical evaluation of both sides of each identity; (c) simulation of the object the
theorem describes, with error bars, compared with the prediction \emph{without fitted constants}. This section
records what was done, what it found, and what remains unverified.

\subsection{What rests on quoted theorems}
\begin{itemize}\setlength\itemsep{2pt}
\item \textbf{Theorem 20.1 of \citet{br}} (Edgeworth expansion for functions of bounded oscillation under Cram\'er's
condition, $s=3$) is used in Lemma~\ref{lem:edge} and Proposition~\ref{prop:lin}, hence in
Theorems~\ref{thm:cf}, \ref{thm:beff}, \ref{thm:state}. \emph{The statement is quoted from secondary sources and from
memory of the monograph; it has not been re-read in the original for this document.} Before submission the exact form
of the oscillation term (the radius $e^{-cb}$ and the uniformity in translations of $f$) should be checked against the
book. The discrete-noise theorems (Section~\ref{sec:discrete}) do not use it and are self-contained.
\item \textbf{Esseen's lattice expansion} \citep[\S43]{gk} is used in Theorem~\ref{thm:sign}(b); the non-lattice case (a)
uses the classical $o(b^{-1/2})$ expansion for non-lattice laws with three moments.
\item \textbf{van der Corput's lemma in several variables} \citep[Ch.~VIII, \S2.2, Prop.~5]{stein} is used in
Lemma~\ref{lem:matrices}; only the dependence of the constant on a lower bound of one derivative and on the $C^{3}$ norm
of the phase is used.
\item \textbf{Hoeffding's comparison} of sampling with and without replacement \citep[Thm.~4]{hoeffding} is used for
(E3) in Proposition~\ref{prop:worclass}.
\item \textbf{Perturbation of polar factors} \citep{lisun} is used in Lemma~\ref{lem:osc}; the divergence formula
for spectral functions \citep{candes} in Theorem~\ref{thm:two}.
\end{itemize}

\subsection{Computer-assisted steps}
Theorem~\ref{thm:a1} for $K\ge3$: the moment polynomials $\E\tr W^k$, $k\le8$ (Wick enumeration), the identity
\eqref{eq:num}, the positivity of the shifted coefficients, and $52$ rational certificates. Two independent
implementations agree (\texttt{note2/a1/cert.py}, \texttt{finite.py}; \texttt{note2/checks/v01\_wick\_moments.py},
\texttt{v02\_a1\_proof.py}); all arithmetic is exact (integers and rationals). In addition $q_{K,d}$ was evaluated to
high precision for all sizes used in the figures and the inequality checked directly.

\subsection{Independent reviews}
\begin{center}\small
\begin{tabular}{@{}p{0.27\textwidth}p{0.20\textwidth}p{0.47\textwidth}@{}}
\toprule
Scope & Report & Outcome\\
\midrule
Note 1 (law under Cram\'er, two-moment theorem, momentum, state dependence) & \texttt{checks/review\_verify.py} &
independent scratch verification of the identities of Theorem~\ref{thm:two} (symbolic and Monte Carlo) and of the checks \texttt{c01}--\texttt{c07}\\
Note 2 (class of sums, codimension theorem, sign theorem, $a_1<0$) & \texttt{note2/referee\_report2.md}, \texttt{note2/referee/} &
tie case of Theorem~\ref{thm:sign} added; additive (not relative) error for weighted sums; a bug in the exact
$2\times2$/$2\times3$ evaluator (rank-deficient atoms) found, fixed, tables regenerated\\
Convergence theorem (Theorem~\ref{thm:conv}) & \texttt{verification/referee\_momentum.md} &
correct; wording fixes (non-anticipating $O_t$, measurability via joint continuity) incorporated\\
Orbital lemma, Laplacian criterion (Thm.~\ref{thm:laplace}, Lem.~\ref{lem:orbit}, \ref{lem:matrices}) &
\texttt{verification/referee\_orbital.md} & correct; all exponents re-derived; span claim verified at every rank for sizes up to
$5\times5$; the mixed derivative $AyB$ shown to be indispensable; sharpness of $\kappa=\frac12$ for $\SO(3)^2$ confirmed numerically\\
Sampling without replacement (Thm.~\ref{thm:wor}) & \texttt{verification/referee\_population.md} & correct, cosmetic fixes only; conditioning identity matched to full subset enumeration ($5\cdot10^{-16}$); exact conditional cumulants confirmed in rational arithmetic; order-$h^{-2}$ term written out; sign reversal and exact zero at $p=\frac12$ reproduced on lattice and non-lattice populations ($\R^3$, Euclidean)\\
\bottomrule
\end{tabular}
\end{center}

\subsection{Numerical checks, by statement}
{\small
\begin{longtable}{@{}p{0.2\textwidth}p{0.33\textwidth}p{0.41\textwidth}@{}}
\toprule
Statement & Check (script) & Result\\
\midrule
\endhead
Thm.~\ref{thm:two}(a) two-constant structure & \texttt{checks/c01} & invariant-theory Gram matrix exact; $T_I$ matches $a_1\cA+a_2\cB$ on random tensors\\
Thm.~\ref{thm:two}(b),(c),(e) & \texttt{checks/c02}; \texttt{papers/experiments} & Stein identity, bounds on $q$, $q_{2,2}=(3\sqrt2-2)/7$; Marchenko--Pastur limit\\
Thm.~\ref{thm:cf}, Cor.~\ref{cor:compare} & \texttt{checks/c03} & five non-product noise laws, isotropic and Kronecker covariance: $b\,G_b\to\gs$\\
Thm.~\ref{thm:beff} & \texttt{checks/c04} & $\beff$ confirmed; the variance-effective rule is off by the predicted factor\\
Thm.~\ref{thm:state} & \texttt{checks/c06}, \texttt{c07}; real-data experiment & zeros at $W_\ast+H^{-1}\gs/b$; plateau $0.3139$ predicted for the finite-sum instance of the original manuscript (measured there: $0.308$ at $b=32$)\\
Thm.~\ref{thm:conv} & \texttt{verification/momentum/} & scalar lattice noise with momentum: limit equals the median of the filtered law; interval of medians gives a random limit; $2\times2$ nuclear case\\
Thm.~\ref{thm:sign} & \texttt{note2/checks/e01} (exact) & $b\,G_b^\st=0,1,-1$ for $b\equiv0,1,2\pmod3$ at every $b$ computed\\
Thm.~\ref{thm:codim} & \texttt{e01} (exact), sizes $1\times2$, $1\times3$, $2\times3$ & convergence to $\gs$ in all residue classes at the stated rates\\
Prop.~\ref{prop:lattice2} & \texttt{note2/checks/e02} & predicted coefficients $-1.012$, $+0.333$, $-0.178$; exact values within a few percent\\
Thm.~\ref{thm:nuclear}, square & \texttt{e01} ($2\times2$, exact); \texttt{verification/orbital/t6} ($3\times3$, Monte Carlo) & convergence in all classes, after transient locking at $0$ in the class $b\equiv0$\\
Thm.~\ref{thm:wor} & \texttt{checks/c05}; \texttt{verification/population/}; figure & exact complement duality on enumerated populations; factor $\frac{n-2b}{n-2}$ and sign change\\
Cor.~\ref{cor:mom} & figure script & iteration with lattice noise and momentum\\
\bottomrule
\end{longtable}}

\subsection{Exact values for lattice noise}
\begin{table}[h]\centering\small
\IfFileExists{../papers/colt/tab_exact.tex}{\begin{tabular}{@{}lccccc@{}}
\toprule
geometry & $\gs$ & $b\equiv0$ & $b\equiv1$ & $b\equiv2$ & largest $b$\\
\midrule
sign ($1\times1$) & $0.1667$ & $0$ & $+1$ & $-1$ & 3000--3002\\
Euclidean $1\times2$ & $0.1250$ & $0.1089$ & $0.1303$ & $0.1222$ & 4096--4098\\
Euclidean $1\times3$ & $0.1000$ & $0.0992$ & $0.1002$ & $0.0999$ & 1000--1002\\
nuclear $2\times3$ & $0.0947$ & $0.0847$ & $0.0960$ & $0.0937$ & 100--102\\
nuclear $2\times2$ & $0.1138$ & $0.0742$ & $0.1176$ & $0.1111$ & 1537--2049\\
nuclear $3\times3$ & $0.0875$ & $0.0871\pm0.0021$ & $0.0866\pm0.0022$ & $0.0883\pm0.0025$ & 3000--3002\\
\bottomrule
\end{tabular}
}{(table pending)}
\caption{Scaled bias $b\,G_b^\st$ for i.i.d.\ entries in $\{-1,2\}$ with probabilities $\frac23,\frac13$ at the largest batch
size computed in each residue class, and the prediction $\gs$.}\label{tab:exact}
\end{table}

\subsection{Known gaps and non-claims}
\begin{itemize}\setlength\itemsep{2pt}
\item Proposition~\ref{prop:lattice2} is a formal derivation. It is not used by any theorem.
\item The exponents $\gamma$ of Theorem~\ref{thm:nuclear} for square matrices are upper bounds on the error; the numerics
suggest they are not sharp. No lower bound is claimed.
\item For general covariance $\Sigma$ the coefficient $\gs$ has no closed form here; Theorem~\ref{thm:two} is for
$\Sigma=\sigma^2I$.
\item Theorem~\ref{thm:state} assumes a uniform expansion (hypothesis (iii)), verified only for linear noise models
(Proposition~\ref{prop:lin}); for finite sums with parameter-dependent discrete noise the location of the fixed point
is supported by experiment, and by Theorem~\ref{thm:nuclear} at each fixed parameter, but not proved uniformly.
\item Theorem~\ref{thm:conv} is for strongly convex smooth $F$ and i.i.d.\ noise independent of the parameter.
\end{itemize}


\bibliographystyle{plainnat}
\bibliography{shared/refs}
\end{document}
